% !TEX TS-program = xelatex
% !TEX encoding = UTF-8

\documentclass{itkmitlproject}

\usepackage{afterpage}
\usepackage{graphicx,amsmath,latexsym,amssymb,amsthm}
\usepackage{indentfirst}
\usepackage{hyperref}
\usepackage{hyphenat}
\usepackage[numbers,square]{natbib}
\usepackage{lipsum}
\usepackage{booktabs}
\usepackage{multirow}
\usepackage{xcolor}
\usepackage{tikz}
\usepackage{float}
\usetikzlibrary{shapes, arrows, positioning, fit, calc, backgrounds}

% \useEnglish
\useThai

\graphicspath{ {images/} }

%1. Your thesis title (THAI)
\newcommand{\ThesisTitle}{การพัฒนาแพลตฟอร์มสำหรับการสนับสนุนการวินิจฉัยรอยโรคทางช่องปากและผสมผสานการใช้งานระบบผู้เชี่ยวชาญสำหรับทันตแพทย์}
%2. Your thesis title (ENG)
\newcommand{\ThesisTitleENG}{Development of an Oral Lesion Diagnosis Support Platform and Integrated Expert System}
%3. Author name
\newcommand{\Author}{นายต้นน้ำ ห้วยหงษ์ทอง}
%4. Author name ENG
\newcommand{\AuthorENG}{Tonnam Huayhongthong}
%5. Author student ID
\newcommand{\SId}{66070259}
%6. Author 2 name
\newcommand{\AuthorTwo}{นายธิติพัทธ์ สนประเสริฐกุล}
%7. Author 2 name ENG
\newcommand{\AuthorTwoENG}{add yo eng name twin, idk how to spell}
%8. Author 2 student ID
\newcommand{\SIdTwo}{66070271}
%9. Advisor name
\newcommand{\Advisor}{ผู้ช่วยศาสตราจารย์ ดร.สุวิทย์ ภูมิฤทธิกุล}
%10. Advisor name ENG
\newcommand{\AdvisorENG}{Asst. Prof. Dr. Suvit Poomrittigul}
%11. Advisor name
\newcommand{\AdvisorTwo}{}
%12. Advisor name ENG
\newcommand{\AdvisorTwoENG}{}
%13. ภาคเรียนที่
\newcommand{\Sem}{1}
%14. ปีการศึกษา พ.ศ.
\newcommand{\AcaY}{2569}
%15. ปีการศึกษา ค.ศ.
%\newcommand{\AcaYAD}{2019}
\newcommand{\AcaYAD}{\the\numexpr\AcaY-543\relax}
%16. วันส่งรายงาน
\newcommand{\SubD}{20 กันยายน พ.ศ. 2569}
%17. Copyright year (AD)
\newcommand{\CopyrightYAD}{2026}

\newcommand{\Qualification}{วิทยาศาสตรบัณฑิต}
\newcommand{\QualificationENG}{Bachelor of Science}
%\newcommand{\Program}{เทคโนโลยีสารสนเทศ}
%\newcommand{\ProgramENG}{Information Technology}
\newcommand{\Program}{วิชาเทคโนโลยีสารสนเทศ}
\newcommand{\ProgramENG}{Information Technology}



\begin{document}
    \frontmatter
    \lhead{}\rhead{}\chead{}\lfoot{}\cfoot{\thepage}\rfoot{}
    \makecover
    \makeinnercover
    \makeengcover
    \makecopyrightcover
    \makethesiscert
    \makeprojectcert
%	\sloppy
    % Setting margin for page numbering on frontmatter
    \newgeometry{top=1in, bottom=1in, left=1.5in, right=1in, includefoot}
    \setcounter{page}{1}
    \pagenumbering{Roman}

    \makeabstract{
    
        ปัจจุบันการตรวจคัดกรองรอยโรคทางช่องปากเพื่อประเมินความเสี่ยงมะเร็งมักประสบปัญหาความล่าช้า เนื่องจากความซับซ้อนของโรคและข้อจำกัดด้านจำนวนทันตแพทย์เฉพาะทาง ที่ผ่านมาได้มีการพัฒนาระบบปัญญาประดิษฐ์เพื่อช่วยคัดกรองรอยโรคในระดับ งานวิจัยต้นแบบ จนประสบความสำเร็จ โครงงานนี้จึงมีแนวคิดในการต่อยอดและยกระดับงานวิจัยดังกล่าวให้เป็นแพลตฟอร์มเว็บแอปพลิเคชันสนับสนุนการตัดสินใจสำหรับทันตแพทย์ที่พร้อมนำไปประยุกต์ใช้งานจริงทางการแพทย์

        ระบบที่พัฒนาขึ้นจะรับข้อมูลภาพถ่ายรอยโรคและประวัติทางคลินิกของผู้ป่วย นำเข้าสู่กระบวนการประมวลผลร่วมกับปัญญาประดิษฐ์เพื่อประเมินผลความเสี่ยงออกมาเป็นคะแนนที่แม่นยำ นอกจากนี้ ยังมีการพัฒนาระบบการให้ข้อมูลย้อนกลับ  บนหน้าแดชบอร์ดที่ช่วยให้ทันตแพทย์สามารถตรวจสอบและแก้ไขผลลัพธ์ของปัญญาประดิษฐ์ เพื่อรวบรวมเป็นชุดข้อมูลคุณภาพสูงสำหรับการพัฒนาความแม่นยำในอนาคต

        การพัฒนาระบบมุ่งเน้นการออกแบบสถาปัตยกรรมคลาวด์ในรูปแบบไมโครเซอร์วิส ควบคู่กับการนำหลักการส่งมอบระบบอัตโนมัติ และ Automated Testing มาประยุกต์ใช้เพื่อรองรับการขยายตัวและเพิ่มความเสถียรในการทำงาน พร้อมทั้งพัฒนาระบบยืนยันตัวตน และการจัดการสิทธิ์ที่รัดกุม เพื่อให้การจัดการข้อมูลผู้ป่วยมีความปลอดภัยสูงสุดตามมาตรฐานการคุ้มครองข้อมูลส่วนบุคคล
    }

    \makeabstracteng{

        Currently, the screening process for oral lesions often faces delays due to the complexity of the diseases and a shortage of specialized dentists. While previous efforts have successfully developed an artificial intelligence screening model at a "Research Prototype" stage, this project aims to extend and elevate this foundation into a Production-Ready web-based decision-support platform for dentists.

        The platform is designed to process oral lesion images alongside patient clinical history, utilizing an integrated AI processing engine to evaluate and generate a reliable clinical risk score. Additionally, the web dashboard features a "Feedback Loop" mechanism that allows dentists to review and manually correct AI predictions, thereby systematically curating high-quality data to improve future model accuracy.

        The system development focuses on implementing a cloud-based Microservices architecture, integrating CI/CD pipelines, and utilizing Automated Testing to ensure deployment stability and scalability. Furthermore, secure authentication and authorization systems are incorporated to guarantee maximum patient data privacy and strict adherence to the Personal Data Protection Act (PDPA).
    }

    \makeack

    \newpage
    \addcontentsline{toc}{chapter}{\contentsname}
    \tableofcontents

    \newpage
    \addcontentsline{toc}{chapter}{\listtablename}
    \listoftables

    \newpage
    \addcontentsline{toc}{chapter}{\listfigurename}
    \listoffigures

    % Reset frontmatter page numbering margin, back to original margin from class file
    \restoregeometry

    \mainmatter
    \lhead{}\rhead{\thepage}\chead{}\lfoot{}\cfoot{}\rfoot{}

    \chapter{บทนำ}
\label{chapter:introduction}

\section{ที่มาและความสำคัญ}

    ในปัจจุบัน เทคโนโลยี AI ได้เข้ามามีบทบาทสำคัญในวงการแพทย์ โดยจากโครงงานการพัฒนาระบบวิเคราะห์รอยโรคทางช่องปากก่อนหน้านี้ ได้ประสบความสำเร็จในการสร้างแบบจำลอง AI ที่สามารถวิเคราะห์โรคได้อย่างแม่นยำ ซึ่งถือเป็นรากฐานที่ดีเยี่ยมในการช่วยทันตแพทย์คัดกรองผู้ป่วย

    อย่างไรก็ตาม ระบบดังกล่าวยังอยู่ในระดับ "งานวิจัยต้นแบบ" เมื่อพิจารณาในมุมมองของวิศวกรรมซอฟต์แวร์ พบว่าระบบเดิมยังมีข้อจำกัดในการนำไปให้บริการจริง เช่น โครงสร้างที่กินทรัพยากรคลาวด์สูงทำให้มีต้นทุนแพง ขาดระบบการจัดการวงจรชีวิตซอฟต์แวร์ ที่ทำให้การอัปเดตระบบทำได้ยาก ตลอดจนยังขาดระบบยืนยันตัวตนและการจัดการข้อมูลที่ปลอดภัยตามมาตรฐาน PDPA

    ด้วยเหตุนี้ โครงงานนี้จึงเกิดขึ้นเพื่อต่อยอดและปิดช่องโหว่ดังกล่าว โดยมุ่งเน้นการปรับปรุงสถาปัตยกรรมระบบใหม่ทั้งหมด นำหลักการ CI/CD มาใช้ ปรับแต่งโมเดล AI ให้ทำงานเร็วและประหยัดต้นทุนขึ้น รวมถึงสร้าง Dashboard เพื่อให้แพทย์สามารถตรวจสอบผลและส่ง Feedback กลับสู่ระบบได้ ซึ่งจะช่วยเปลี่ยนโปรเจกต์นี้ให้กลายเป็นซอฟต์แวร์ทางการแพทย์ที่เสถียร ปลอดภัย และพร้อมใช้งานจริงอย่างยั่งยืน

\section{วัตถุประสงค์}

    1.2.1 เพื่อพัฒนาระบบผู้เชี่ยวชาญวิเคราะห์รอยโรคทางช่องปากให้เป็นแพลตฟอร์มที่เสถียรและพร้อมใช้งานจริง โดยประยุกต์ใช้วงจรชีวิตการพัฒนาซอฟต์แวร์ และการนำส่งระบบอัตโนมัติ (CI/CD)
    
    1.2.2 เพื่อสร้างสภาพแวดล้อมที่ปลอดภัยสำหรับการประมวลผลข้อมูลรอยโรคผู้ป่วย โดยพัฒนาระบบรักษาความปลอดภัยและการยืนยันตัวตนให้สอดคล้องกับมาตรฐานการคุ้มครองข้อมูลส่วนบุคคล (PDPA)
    
    1.2.3 เพื่อเพิ่มประสิทธิภาพและลดต้นทุน การทำงานของแบบจำลองปัญญาประดิษฐ์วิเคราะห์รอยโรคเดิม ผ่านเทคนิค Model Optimization และการปรับสถาปัตยกรรมคลาวด์ให้รองรับผู้ใช้งานจำนวนมากแบบคู่ขนาน
    
    1.2.4 เพื่อสร้างเว็บแอปพลิเคชันสนับสนุนการตัดสินใจสำหรับทันตแพทย์ ที่อำนวยความสะดวกในการตรวจสอบรอยโรค พร้อมระบบให้ข้อมูลย้อนกลับ (Feedback Loop) สำหรับรวบรวมข้อมูลไปพัฒนาความแม่นยำของโมเดลในอนาคต

\section{วิธีการดำเนินงาน}

    1.3.1 ศึกษาและวิเคราะห์โครงสร้างระบบเดิม และ รวบรวมข้อมูลภาพถ่ายช่องปาก ข้อมูลจากการซักประวัติและผลวินิจฉัยจากผู้เชี่ยวชาญ

    1.3.2 ออกแบบสถาปัตยกรรมระบบใหม่ (Microservices) เพื่อรองรับการทำงานที่เสถียรยิ่งขึ้น
    
    1.3.3 นำแบบจำลองเดิมมาทำการปรับแต่ง (Model Optimization) ให้อยู่ในรูปแบบที่เหมาะสม เพื่อลดขนาดและเพิ่มความเร็วในการอนุมานผล

    1.3.4 พัฒนา Web Application และ Doctor Dashboard สำหรับทันตแพทย์โดยให้ทันตแพทย์สามารถกดยืนยันหรือแก้ไขผลการทำนายของ AI ได้เพื่อนำข้อมูลไปพัฒนาความแม่นยำต่อไปในอนาคต

    1.3.5 จัดทำ Automated Testing สำหรับตรวจสอบ API และทดสอบความปลอดภัยของระบบ (Security/Vulnerability Scanning) รวมถึงการทำ Load Testing

    1.3.6 ออกแบบและติดตั้งระบบ CI/CD Pipeline ควบคู่กับการติดตั้งระบบ Centralized Logging เพื่อตรวจจับข้อผิดพลาด (Error Tracking) ในการนำส่งและดูแลรักษาระบบ

    1.3.7 สรุปผลการดำเนินงานของระบบ

\section{ขอบเขตโครงงาน}

    1.4.1  พัฒนาส่วนติดต่อผู้ใช้งานสำหรับทันตแพทย์ เพื่อใช้ในการกรอกข้อมูลซักประวัติ อัปโหลดรูปภาพรอยโรค และแสดงผลการวิเคราะห์ความเสี่ยงจากแบบจำลอง AI 
    
    1.4.2 ออกแบบและพัฒนาระบบหลังบ้าน ในรูปแบบ Microservices บนสภาพแวดล้อม Cloud เพื่อให้รองรับการใช้งานพร้อมกันของคลินิกหรือทันตแพทย์หลายราย

    1.4.3 พัฒนาระบบยืนยันตัวตน ระบบจัดการสิทธิ์ผู้ใช้งาน และการเข้ารหัสข้อมูล เพื่อให้การจัดเก็บข้อมูลผู้ป่วยสอดคล้องกับพ.ร.บ. คุ้มครองข้อมูลส่วนบุคคล (PDPA)

    1.4.4 พัฒนาระบบที่อนุญาตให้ทันตแพทย์สามารถกดยืนยัน หรือแก้ไขผลการวินิจฉัยของ AI ได้ เพื่อรวบรวมข้อมูล Ground Truth กลับไปใช้พัฒนาความแม่นยำของโมเดลในอนาคต

    1.4.5 ติดตั้งระบบ Continuous Integration / Continuous Deployment ควบคู่กับการทำ Automated Testing เพื่อให้การอัปเดตระบบเป็นไปอย่างรวดเร็วและปลอดภัย

    1.4.6  ติดตั้งระบบ Centralized Logging เพื่อตรวจสอบการทำงานของระบบและแจ้งเตือนข้อผิดพลาด แบบเรียลไทม์

\section{ประโยชน์ที่คาดว่าจะได้รับ}

    1.5.1 สามารถยกระดับจากงานวิจัยต้นแบบ ให้กลายเป็นแพลตฟอร์มที่เสถียร รองรับผู้ใช้งานจำนวนมาก และสามารถนำไปติดตั้งใช้งานจริงในคลินิกหรือโรงพยาบาลได้

    1.5.2 สถาปัตยกรรมระบบและโมเดลที่ได้รับการ Optimize จะช่วยลดค่าใช้จ่ายในการเช่าทรัพยากร Cloud Server ลงได้อย่างมีนัยสำคัญ

    1.5.3 ผู้ใช้งานและผู้ป่วยสามารถมั่นใจได้ว่าข้อมูลสุขภาพส่วนบุคคลจะถูกจัดเก็บและประมวลผลอย่างปลอดภัยตามมาตรฐาน PDPA

    1.5.4 ช่วยเพิ่มประสิทธิภาพในการบริหารจัดการเอกสารและการสร้างผลลัพธ์ของงานเมื่อผ่านการอนุมัติครบตามกระบวนการ

    1.5.5 การมีระบบ Feedback Loop จะช่วยให้ทีมผู้พัฒนาโมเดลมีชุดข้อมูลที่ผ่านการคัดกรองจากผู้เชี่ยวชาญ (แพทย์) เพื่อนำไป Train โมเดลให้แม่นยำยิ่งขึ้นในเวอร์ชันถัดไป

    1.5.6 ช่วยลดระยะเวลาในการวิเคราะห์เบื้องต้น ทำให้ทันตแพทย์สามารถคัดกรองรอยโรคที่อาจนำไปสู่โรคมะเร็งช่องปากได้รวดเร็วและเป็นระบบมากขึ้น

    1.5.7 ได้องค์ความรู้และประสบการณ์ในการออกแบบและพัฒนาระบบที่ผสานการจัดการข้อมูล ระบบ Workflow ระบบอนุมัติ และ AI เข้าด้วยกัน

    1.5.8 ได้ระบบต้นแบบที่สามารถต่อยอดหรือปรับใช้กับกระบวนการงานอื่นภายในสถาบันในอนาคตได้

    % ==============================================================================
% บทที่ 2: การทบทวนวรรณกรรมที่เกี่ยวข้อง (Literature Review)
% โครงการ: การพัฒนาแพลตฟอร์มสำหรับการสนับสนุนการวินิจฉัยรอยโรคทางช่องปาก
%          และผสมผสานการใช้งานระบบผู้เชี่ยวชาญสำหรับทันตแพทย์
% ==============================================================================

\chapter{การทบทวนวรรณกรรมที่เกี่ยวข้อง}
\label{chapter:literature-review}

การตรวจคัดกรองและวินิจฉัยรอยโรคในช่องปาก (Oral Lesion Screening and Diagnosis) โดยเฉพาะการประเมินความเสี่ยงต่อการเกิดโรคมะเร็งช่องปาก (Oral Cancer) ถือเป็นขั้นตอนที่มีความสำคัญอย่างยิ่งต่อการเพิ่มอัตราการรอดชีวิตของผู้ป่วย โดยมะเร็งช่องปากส่วนใหญ่จัดอยู่ในกลุ่มมะเร็งเซลล์สความัส (Oral Squamous Cell Carcinoma: OSCC) ซึ่งมักมีจุดเริ่มต้นมาจากการเปลี่ยนแปลงของเนื้อเยื่อบุผิวที่เป็นรอยโรคก่อนมะเร็ง (Oral Potentially Malignant Disorders: OPMDs) อย่างไรก็ตาม ในทางคลินิกพบว่ารอยโรคระยะแรกเริ่มมักมีลักษณะทางสัณฐานวิทยาใกล้เคียงกับรอยโรคทั่วไปที่ไม่ร้ายแรง ทำให้การแยกแยะโรคต้องพึ่งพาความชำนาญของทันตแพทย์เฉพาะทางด้านเวชศาสตร์ช่องปาก ซึ่งยังมีจำนวนจำกัดและกระจุกตัวอยู่ในโรงพยาบาลระดับศูนย์

จากการศึกษาและต่อยอดโครงงานวิจัยของรุ่นพี่ในปีก่อนหน้าทั้ง 2 รุ่น พบว่า:
\begin{enumerate}
    \item \textbf{โครงงานวิจัยของรุ่นพี่กลุ่มที่ 1 (ปีการศึกษา 2567):} ได้พัฒนาระบบตรวจจับรอยโรคโดยใช้แบบจำลองตรวจจับวัตถุ (Object Detection) ตระกูล YOLO ร่วมกับการประมวลผลข้อมูลอาการของผู้ป่วยด้วยกฎเงื่อนไข (Rule-based) แม้จะสามารถตีกรอบ Bounding Box ระบุตำแหน่งรอยโรคได้ดี แต่การตัดสินใจยังเป็นการวิเคราะห์แบบแยกส่วน และมีข้อจำกัดอย่างมากเมื่อภาพถ่ายไม่สามารถตรวจจับรอยโรคได้ (No Detection) หรือเมื่อผู้ป่วยมีอาการที่ขัดแย้งกัน
    \item \textbf{โครงงานวิจัยของรุ่นพี่กลุ่มที่ 2 (ปีการศึกษา 2568):} ได้พัฒนาอัลกอริทึมการเรียนรู้เชิงลึกแบบหลายโมดาลิตี (Multimodal Deep Learning) โดยใช้โครงข่าย DenseNet-121 วิเคราะห์ภาพถ่ายร่วมกับแบบจำลองภาษาไทย WangchanBERTa วิเคราะห์ข้อความซักประวัติ ซึ่งให้ค่าความแม่นยำสูงถึงร้อยละ 95.5 และทนทานต่อข้อมูลอาการที่ขัดแย้งกันได้ดี แต่มีข้อจำกัดคือเป็นแบบจำลองที่ทำนายเฉพาะภาพรวม (Whole-image classification) ไม่สามารถตีกรอบระบุพิกัดตำแหน่งรอยโรคบนภาพได้
\end{enumerate}

แม้ว่าระบบต้นแบบของรุ่นพี่ทั้งสองกลุ่มจะประสบความสำเร็จอย่างสูงในระดับห้องปฏิบัติการ แต่เมื่อพิจารณาในมุมมองของวิศวกรรมซอฟต์แวร์และการนำไปให้บริการจริงทางการแพทย์ พบว่าระบบเดิมยังคงอยู่ในระดับ ``งานวิจัยต้นแบบ (Research Prototype)'' ที่มีข้อจำกัดด้านโครงสร้างระบบ เช่น สถาปัตยกรรมที่กินทรัพยากรประมวลผลคลาวด์สูง ทำให้มีต้นทุนค่าบริการแพง ขาดระบบการจัดการวงจรชีวิตซอฟต์แวร์และการส่งมอบระบบอัตโนมัติ (CI/CD) ขาดระบบการทดสอบอัตโนมัติ (Automated Testing) ตลอดจนยังไม่มีระบบยืนยันตัวตนและการจัดการสิทธิ์ที่ปลอดภัยตามมาตรฐานพระราชบัญญัติคุ้มครองข้อมูลส่วนบุคคล (PDPA) และยังขาดหน้าต่างการทำงานสำหรับทันตแพทย์ (Doctor Dashboard) ที่มีระบบวงจรข้อมูลย้อนกลับ (Feedback Loop) สำหรับให้แพทย์ตรวจสอบ ยืนยัน หรือแก้ไขผลลัพธ์ของ AI เพื่อเก็บรวบรวมเป็นชุดข้อมูลคุณภาพสูงไปพัฒนาแบบจำลองให้แม่นยำยิ่งขึ้นในอนาคต

ดังนั้น เพื่อเป็นพื้นฐานในการออกแบบและพัฒนาแพลตฟอร์มสนับสนุนการตัดสินใจสำหรับทันตแพทย์ที่พร้อมใช้งานจริง (Production-Ready Platform) และบูรณาการระบบผู้เชี่ยวชาญทั้งสองแนวทางเข้าด้วยกัน เนื้อหาในบทนี้จึงมุ่งเน้นการทบทวนวรรณกรรม สถาปัตยกรรมระบบ วิศวกรรมซอฟต์แวร์ และเทคโนโลยีที่เกี่ยวข้อง โดยแบ่งออกเป็น 7 หัวข้อหลัก ได้แก่:
\begin{enumerate}
    \item สถาปัตยกรรมซอฟต์แวร์และระบบคลาวด์ไมโครเซอร์วิส
    \item เทคโนโลยีส่วนต่อประสานผู้ใช้ แดชบอร์ดแพทย์ และการสื่อสาร
    \item วิศวกรรมซอฟต์แวร์ การทดสอบอัตโนมัติ และความปลอดภัยของข้อมูลสุขภาพ
    \item ภาพรวมสถาปัตยกรรมปัญญาประดิษฐ์และมาตรวัดประสิทธิภาพ
    \item การทบทวนงานวิจัยและปริญญานิพนธ์ที่เกี่ยวข้อง
    \item สรุปแนวทางการออกแบบสถาปัตยกรรมและกรอบแนวคิดของโครงงาน
\end{enumerate}

% ------------------------------------------------------------------------------
% 2.1 แนวคิดทางการแพทย์และพยาธิวิทยาของรอยโรคในช่องปาก
% ------------------------------------------------------------------------------
\section{แนวคิดทางการแพทย์และพยาธิวิทยาของรอยโรคในช่องปาก}
\label{sec:medical_pathology}

\subsection{รอยโรคเป้าหมาย 5 กลุ่ม}
ในการพัฒนาระบบคัดกรอง โครงงานนี้ครอบคลุมกลุ่มรอยโรคสีขาวและแผลในช่องปากที่มีความสำคัญทางคลินิก 5 ประเภท ได้แก่:
\begin{enumerate}
    \item \textbf{แผลเปื่อยหรือแผลร้อนใน (Aphthous Ulcer / Ulcer):} เป็นแผลหลุมตื้น ขอบแผลแดง ตรงกลางเป็นสีขาวเหลือง เกิดจากการอักเสบ ผู้ป่วยมีอาการเจ็บปวดชัดเจนเมื่อสัมผัส มักหายได้เองภายใน 1--2 สัปดาห์ แต่หากแผลเรื้อรังเกิน 3 สัปดาห์โดยไม่หายหรือมีขอบแข็งยกตัวขึ้น จะต้องสงสัยมะเร็งช่องปาก
    \item \textbf{โรคติดเชื้อราแคนดิดาในช่องปาก (Oral Candidiasis):} เกิดจากเชื้อรา \textit{Candida albicans} มีลักษณะเป็นแผ่นฝ้าขาวคล้ายคราบนม \textbf{ลักษณะเด่นคือสามารถใช้ผ้ากอซขูดหรือเช็ดออกได้ (Scrapable plaque)} และเมื่อเช็ดออกจะพบผิวเยื่อบุสีแดงด้านล่าง
    \item \textbf{เส้นขาวแนวกระพุ้งแก้ม (Linea Alba Buccalis):} เป็นการเปลี่ยนแปลงทางกายวิภาคปกติ (Normal variation) เกิดจากแรงกดหรือแรงเสียดสีซ้ำๆ จากแนวสบฟัน เป็นเส้นขาวนูนตามแนวนอนของกระพุ้งแก้ม \textbf{ไม่มีอาการเจ็บหรือแสบใดๆ และขูดไม่ออก} ไม่มีความเสี่ยงต่อการเกิดมะเร็ง
    \item \textbf{รอยโรคแผ่นขาว (Oral Leukoplakia):} เป็นรอยโรคก่อนมะเร็ง (OPMD) ที่สำคัญที่สุด มีลักษณะเป็นแผ่นหรือปื้นสีขาวหนาตัว \textbf{ขูดไม่ออก} สัมพันธ์กับพฤติกรรมการสูบบุหรี่ ดื่มสุรา และเคี้ยวหมากเรื้อรัง มีอัตราการกลายพันธุ์เป็นมะเร็งช่องปากประมาณร้อยละ 1--9
    \item \textbf{โรคไลเคน พลานัสในช่องปาก (Oral Lichen Planus: OLP):} เป็นโรคอักเสบเรื้อรังจากระบบภูมิคุ้มกัน มีลักษณะเป็นเส้นสีขาวสานกันคล้ายร่างแห (Wickham's striae) หรือแผลแดงถลอก ผู้ป่วยมักมีอาการแสบร้อนเมื่อรับประทานอาหารเผ็ด จัดเป็นรอยโรคกลุ่ม OPMD ที่มีโอกาสกลายเป็นมะเร็งได้ประมาณร้อยละ 0.5--2
\end{enumerate}

\begin{table}[htbp]
\centering
\small
\caption{สรุปการเปรียบเทียบลักษณะทางคลินิกและพยาธิวิทยาของรอยโรคเป้าหมายทั้ง 5 ชนิด}
\label{tab:clinical_pathology_summary}
\begin{tabular}{p{2.8cm}p{3.2cm}p{2.6cm}p{4.2cm}}
\toprule
\textbf{ชื่อรอยโรค} & \textbf{ลักษณะทางกายภาพ} & \textbf{การขูดลอก} & \textbf{อาการทางคลินิกและพฤติกรรมเสี่ยง} \\
\midrule
Ulcer (แผลเปื่อย) & แผลหลุมตื้น ตรงกลางขาวเหลือง ขอบแดง & ขูดไม่ออก & เจ็บปวดมากเวลากินอาหารหรือสัมผัส หายใน 2 สัปดาห์ \\
Candidiasis (เชื้อรา) & แผ่นฝ้าขาวคล้ายคราบนมกระจายตัว & เช็ด/ขูดลอกได้ & แสบปาก ผิวใต้รอยแดง อาจพบจุดเลือดซึม \\
Linea Alba & เส้นขาวนูนทอดตามแนวสบฟัน & ขูดไม่ออก & ไม่มีอาการใดๆ, สัมพันธ์กับแรงสบฟัน \\
Leukoplakia (แผ่นขาว) & แผ่นหรือปื้นขาวขอบเขตชัดเจน & ขูดไม่ออก & สัมพันธ์กับสูบบุหรี่ ดื่มสุรา เคี้ยวหมาก (เสี่ยงมะเร็ง) \\
Lichen Planus & เส้นขาวร่างแห หรือแผลแดงถลอก & ขูดไม่ออก & ปวดแสบปวดร้อนเมื่อกินอาหารเผ็ด (เสี่ยงมะเร็ง) \\
\bottomrule
\end{tabular}
\end{table}

% ------------------------------------------------------------------------------
% 2.2 สถาปัตยกรรมซอฟต์แวร์และระบบคลาวด์ไมโครเซอร์วิส
% ------------------------------------------------------------------------------
\section{สถาปัตยกรรมซอฟต์แวร์และระบบคลาวด์ไมโครเซอร์วิส}
\label{sec:software_architecture}

\subsection{สถาปัตยกรรมแบบ Monolith เทียบกับสถาปัตยกรรม Microservice}
ในระบบต้นแบบเดิมของงานวิจัยรุ่นก่อนหน้า มักใช้สถาปัตยกรรมแบบ Monolith ซึ่งรวมทั้งส่วนติดต่อผู้ใช้ การประมวลผลทางธุรกิจ และการโหลดโมเดลปัญญาประดิษฐ์เข้าไว้ในแอปพลิเคชันเดียวกัน ส่งผลให้เกิดข้อจำกัดหลายประการ:
\begin{itemize}
    \item \textbf{การใช้ทรัพยากรที่ไม่คุ้มค่า:} การรันโมเดล AI ขนาดใหญ่ (เช่น PyTorch, CUDA, Transformers) ภายในเว็บเซิร์ฟเวอร์หลัก ทำให้ตัวเซิร์ฟเวอร์ต้องใช้หน่วยความจำ (RAM) และ CPU สูงตลอดเวลา แม้จะไม่มีการส่งภาพเข้ามาวิเคราะห์
    \item \textbf{ความยากในการขยายขนาด (Scalability):} เมื่อมีผู้ใช้งานเข้าดูเว็บไซต์พร้อมกันจำนวนมาก ระบบไม่สามารถขยายเฉพาะเว็บเซิร์ฟเวอร์ได้ แต่ต้องคัดลอกทั้งอินสแตนซ์ที่รวมโมเดล AI ไปด้วย ซึ่งทำให้ค่าใช้จ่ายบนคลาวด์เพิ่มสูงขึ้นอย่างรวดเร็ว
    \item \textbf{ความเสี่ยงของระบบล่ม (Single Point of Failure):} หากกระบวนการอนุมานผลของ AI เกิดข้อผิดพลาดจากหน่วยความจำไม่เพียงพอ (Out of Memory: OOM) จะส่งผลให้ทั้งเว็บเซิร์ฟเวอร์หยุดทำงานทันที
\end{itemize}

ด้วยเหตุนี้ โครงงานนี้จึงปรับเปลี่ยนสถาปัตยกรรมไปสู่ สถาปัตยกรรมไมโครเซอร์วิส (Microservices Architecture) โดยแยกส่วนการทำงานออกเป็น 2 เซอร์วิสหลัก:
\begin{enumerate}
    \item \textbf{API Gateway Service:} รับคำขอจากผู้ใช้งาน ตรวจสอบสิทธิ์ ตรวจสอบความถูกต้องของข้อมูล และส่งต่องานไปยัง AI Engine
    \item \textbf{AI Engine Inference Service:} ทำหน้าที่โหลดโมเดลเข้าสู่หน่วยความจำและรันการอนุมานผลภาพและข้อความประวัติผู้ป่วยโดยเฉพาะ
\end{enumerate}

\subsection{รูปแบบประตูเชื่อมต่อระบบ (API Gateway Pattern) และ FastAPI}
FastAPI ถูกเลือกใช้เป็นเฟรมเวิร์กหลักในการสร้าง API Gateway และ AI Engine ด้วยคุณสมบัติที่โดดเด่น:
\begin{itemize}
    \item \textbf{Asynchronous Programming (ASGI):} ทำงานบนสถาปัตยกรรมแบบ Asynchronous ควบคู่กับเซิร์ฟเวอร์ Uvicorn ทำให้สามารถรองรับคำขอ (Concurrency) ได้พร้อมกันจำนวนมหาศาลโดยไม่เกิดการบล็อก (Non-blocking I/O)
    \item \textbf{Data Validation \& Serialization ด้วย Pydantic:} ทำหน้าที่ตรวจสอบโครงสร้างข้อมูลอินพุต (เช่น ขนาดและรูปแบบไฟล์ภาพ, ข้อความประวัติ) และจัดรูปแบบผลลัพธ์การประเมิน (Unified Assessment Schema) ให้เป็นมาตรฐานเดียวกันก่อนส่งกลับไปยังผู้ใช้
    \item \textbf{Auto-Generated Documentation:} สร้างเอกสารแนะนำการใช้งาน API อัตโนมัติในรูปแบบ OpenAPI (Swagger UI) ช่วยให้การพัฒนาและเชื่อมต่อกับระบบภายนอกเป็นไปอย่างมีมาตรฐาน
\end{itemize}

\subsection{เทคโนโลยีคอนเทนเนอร์ (Containerization: Docker \& Docker Compose)}
Docker ถูกนำมาใช้ในการบรรจุแอปพลิเคชัน โค้ด รันไทม์ และชุดไลบรารีที่จำเป็นทั้งหมด (เช่น PyTorch, OpenCV, Transformers) เข้าเป็นอิมเมจมาตรฐาน (Container Images) เพื่อตัดปัญหา ``ทำงานได้บนเครื่องผู้พัฒนาแต่รันบนเซิร์ฟเวอร์ไม่ได้ (It works on my machine problem)'' นอกจากนี้ยังใช้ Docker Compose ในการจัดการประสานงานระหว่างคอนเทนเนอร์ Gateway และคอนเทนเนอร์ AI Engine ทำให้สามารถจำลองสภาพแวดล้อมเสมือนจริงได้อย่างสมบูรณ์ทั้งในเครื่องทดสอบและบนคลาวด์


% ------------------------------------------------------------------------------
% 2.3 เทคโนโลยีส่วนต่อประสานผู้ใช้ แดชบอร์ดแพทย์ และการสื่อสาร
% ------------------------------------------------------------------------------
\section{เทคโนโลยีส่วนต่อประสานผู้ใช้ แดชบอร์ดแพทย์ และการสื่อสาร}
\label{sec:ui_dashboard_comm}

\subsection{การพัฒนาเว็บแอปพลิเคชันส่วนติดต่อผู้ใช้ (Web Application Frontend)}
ส่วนติดต่อผู้ใช้ถูกออกแบบเพื่อรองรับการใช้งานที่สะดวกและรวดเร็วสำหรับทันตแพทย์ โดยใช้เทคโนโลยีเว็บสมัยใหม่:
\begin{itemize}
    \item \textbf{Single-Page Application (SPA):} ช่วยให้การตอบสนองต่อการคลิกและเปลี่ยนหน้าเป็นไปอย่างรวดเร็วโดยไม่ต้องรีโหลดหน้าเว็บใหม่ทั้งหมด
    \item \textbf{Tailwind CSS:} เฟรมเวิร์กการจัดแต่งรูปแบบการแสดงผลแบบ Utility-first ที่ช่วยให้การจัดวางองค์ประกอบมีความยืดหยุ่นสูง สวยงาม รองรับการแสดงผลทุกขนาดหน้าจอ (Responsive Web Design) ทั้งบนคอมพิวเตอร์ แท็บเล็ต และสมาร์ตโฟน
    \item \textbf{Form \& File Upload Handling:} มีระบบตรวจสอบความสมบูรณ์ของภาพถ่ายและแบบฟอร์มซักประวัติก่อนส่งข้อมูล เพื่อลดความผิดพลาดที่อาจเกิดขึ้นจากฝั่งผู้ใช้งาน
\end{itemize}

\subsection{แดชบอร์ดสนับสนุนการตัดสินใจของทันตแพทย์และระบบวงจรข้อมูลย้อนกลับ (Doctor Clinical Dashboard \& Feedback Loop)}
หนึ่งในนวัตกรรมสำคัญของแพลตฟอร์มนี้คือ **ระบบวงจรข้อมูลย้อนกลับ (Doctor Feedback Loop)** ซึ่งถูกสร้างขึ้นบนแดชบอร์ดสำหรับทันตแพทย์โดยเฉพาะ:
\begin{enumerate}
    \item \textbf{Clinical Decision Support:} ระบบจะแสดงผลการวิเคราะห์ภาพถ่ายคู่กับแผนที่ความร้อน Grad-CAM พร้อมแสดงกรอบ Bounding Box และคะแนนระดับความเสี่ยงมะเร็ง (Risk Score) เพื่อให้ทันตแพทย์ใช้ประกอบการวินิจฉัย
    \item \textbf{Doctor Verification \& Relabeling:} ทันตแพทย์ผู้เชี่ยวชาญสามารถกดยืนยันผลการทำนายของ AI หรือทำการแก้ไขชนิดรอยโรค (Relabeling) ได้หากพบว่าผลการทำนายของ AI ยังไม่ถูกต้อง
    \item \textbf{Ground Truth Data Curation:} ข้อมูลที่ได้รับการยืนยันหรือแก้ไขจากแพทย์จะถูกจัดเก็บเข้าสู่ฐานข้อมูลคุณภาพสูง (Curated Ground Truth Dataset) โดยอัตโนมัติ ซึ่งทำหน้าที่เป็นขุมข้อมูลสำคัญสำหรับการนำไปฝึกฝนและปรับปรุงโมเดล AI ในรอบถัดไป (Continuous Learning) ทำให้ระบบมีความแม่นยำเพิ่มขึ้นอย่างยั่งยืน
\end{enumerate}

% ------------------------------------------------------------------------------
% 2.4 วิศวกรรมซอฟต์แวร์ การทดสอบอัตโนมัติ และความปลอดภัยของข้อมูลสุขภาพ
% ------------------------------------------------------------------------------
\section{วิศวกรรมซอฟต์แวร์ การทดสอบอัตโนมัติ และความปลอดภัยของข้อมูลสุขภาพ}
\label{sec:software_engineering_pdpa}

\subsection{การผสานรวมและการส่งมอบระบบอย่างต่อเนื่อง (CI/CD Pipelines)}
เพื่อให้การดูแลรักษาและการพัฒนาฟีเจอร์ใหม่เป็นไปอย่างรวดเร็วและปลอดภัย โครงงานนี้นำแนวคิด DevOps และกระบวนการ **CI/CD (Continuous Integration / Continuous Deployment)** มาประยุกต์ใช้ผ่าน GitHub Actions:
\begin{itemize}
    \item \textbf{Continuous Integration (CI):} ทุกครั้งที่มีการแก้ไขและพุชโค้ดเข้าสู่คลังจัดเก็บ ระบบจะทำการรันการตรวจสอบรูปแบบโค้ด (Linting) และรันชุดทดสอบอัตโนมัติ (Automated Testing) ทันที
    \item \textbf{Continuous Deployment (CD):} เมื่อผ่านการทดสอบครบถ้วน ระบบจะทำการสร้าง Docker Image และส่งมอบไปยังสภาพแวดล้อมคลาวด์ (Cloud Run) โดยอัตโนมัติ ช่วยลดความผิดพลาดจากการติดตั้งแบบ Manual
\end{itemize}

\subsection{การทดสอบระบบอัตโนมัติ (Automated Testing)}
การสร้างความเชื่อถือได้ของซอฟต์แวร์ทางการแพทย์จำเป็นต้องมีระบบการทดสอบที่ครอบคลุม:
\begin{enumerate}
    \item \textbf{Unit Testing \& Integration Testing:} ทดสอบฟังก์ชันการทำงานระดับย่อยและการเชื่อมต่อระหว่าง API Gateway และ AI Engine โดยใช้ Pytest และ HTTPX Mocking
    \item \textbf{Load \& Stress Testing:} การจำลองการส่งคำขอเข้ามาพร้อมกันจำนวนมากเพื่อวัดระยะเวลาการตอบสนอง (Response Time / Latency) และตรวจสอบขีดจำกัดของระบบในการรองรับการใช้งานพร้อมกันของคลินิกหลายแห่ง
    \item \textbf{Security \& Vulnerability Scanning:} ตรวจสอบช่องโหว่ความปลอดภัยของไลบรารีและแพ็กเกจที่ใช้งาน เพื่อป้องกันการถูกโจมตีทางไซเบอร์
\end{enumerate}

\subsection{ระบบบันทึกประวัติและการตรวจจับข้อผิดพลาดส่วนกลาง (Centralized Logging)}
ระบบติดตั้งกลไก Centralized Logging เพื่อรวบรวมบันทึกเหตุการณ์ (Logs) จากทั้ง Gateway และ AI Engine เข้าสู่ศูนย์กลาง ช่วยให้ผู้ดูแลระบบสามารถติดตามข้อผิดพลาด (Error Tracking), ตรวจสอบความผิดปกติของการประมวลผลโมเดล และวิเคราะห์สาเหตุของปัญหาได้อย่างรวดเร็วแบบเรียลไทม์

\subsection{ความปลอดภัยและการคุ้มครองข้อมูลส่วนบุคคล (PDPA Compliance for Patient Records)}
ข้อมูลภาพถ่ายช่องปากและประวัติการรักษาจัดเป็นข้อมูลส่วนบุคคลที่มีความอ่อนไหวสูง (Sensitive Personal Data) ระบบจึงได้รับการออกแบบให้สอดคล้องตาม **พระราชบัญญัติคุ้มครองข้อมูลส่วนบุคคล พ.ศ. 2562 (PDPA)**:
\begin{itemize}
    \item \textbf{Authentication \& Role-Based Access Control (RBAC):} มีระบบยืนยันตัวตนของผู้ใช้งาน และแบ่งระดับสิทธิ์การเข้าถึงอย่างชัดเจน เช่น สิทธิ์สำหรับทันตแพทย์ทั่วไปในการส่งตรวจ และสิทธิ์สำหรับทันตแพทย์ผู้เชี่ยวชาญในการเข้าถึงหน้าแดชบอร์ดตรวจสอบและแก้ไขผล
    \item \textbf{Data Anonymization \& Pseudonymization:} ข้อมูลภาพถ่ายและประวัติที่นำมาประมวลผลในระบบ AI จะถูกตัดข้อมูลที่สามารถระบุตัวตนบุคคลได้ (เช่น ชื่อ-นามสกุล หมายเลขบัตรประชาชน) ออกจากชุดข้อมูลประมวลผล
    \item \textbf{Data Encryption:} ข้อมูลทั้งหมดจะถูกเข้ารหัสขณะส่งผ่านเครือข่ายด้วยโปรโตคอล HTTPS (TLS 1.3) และเข้ารหัสข้อมูลขณะจัดเก็บในฐานข้อมูล (Encryption at Rest)
\end{itemize}

% ------------------------------------------------------------------------------
% 2.5 ภาพรวมสถาปัตยกรรมปัญญาประดิษฐ์และมาตรวัดประสิทธิภาพ
% ------------------------------------------------------------------------------
\section{ภาพรวมสถาปัตยกรรมปัญญาประดิษฐ์และมาตรวัดประสิทธิภาพ}
\label{sec:ai_architecture_overview}

\subsection{ภาพรวมแบบจำลองตรวจจับวัตถุ YOLO (You Only Look Once)}
สถาปัตยกรรม YOLO (YOLOv7 และ YOLOv8) ถูกใช้ในระบบเพื่อทำหน้าที่ตรวจจับตำแหน่งและตีกรอบ Bounding Box ระบุพิกัดรอยโรคบนภาพถ่าย โดยแบบจำลองจะรับภาพถ่ายช่องปากเข้ามา สกัดคุณลักษณะหลายระดับผ่าน Backbone และ Neck จากนั้น Detection Head จะทำนายพิกัดกรอบ $[x, y, w, h]$ และค่าความมั่นใจ (Confidence Score) พร้อมทั้งใช้กระบวนการ Non-Maximum Suppression (NMS) เพื่อคัดกรองกรอบที่ซ้ำซ้อนกันออกไป ทำให้แพทย์เห็นตำแหน่งที่ชัดเจนของแผลบนภาพ

\subsection{ภาพรวมแบบจำลองการเรียนรู้เชิงลึกแบบหลายโมดาลิตี (Multimodal AI)}
เพื่อประเมินความเสี่ยงต่อการเกิดมะเร็งช่องปากอย่างแม่นยำ ระบบได้ผสานข้อมูลสองรูปแบบ:
\begin{enumerate}
    \item \textbf{โมดาลิตีภาพถ่าย (DenseNet-121):} ทำหน้าที่สกัดคุณลักษณะทางกายภาพของรอยโรค เช่น สี ขอบเขต และพื้นผิว ผ่าน Dense Block ที่มีการเชื่อมต่อแบบหนาแน่น ส่งผลให้ได้เวกเตอร์คุณลักษณะขนาด $1 \times 1024$ มิติ
    \item \textbf{โมดาลิตีข้อความประวัติผู้ป่วย (WangchanBERTa):} ทำหน้าที่ประมวลผลข้อความอาการภาษาไทย เช่น ประวัติการสูบบุหรี่ การดื่มสุรา หรืออาการเจ็บแสบ ผ่านกลไก Self-Attention ส่งผลให้ได้เวกเตอร์ตัวแทนบริบทภาษาไทยขนาด $1 \times 768$ มิติ
    \item \textbf{Late Decision Fusion:} นำคะแนนผลลัพธ์ของโมเดลภาพและโมเดลข้อความมารวมกันด้วยการถ่วงน้ำหนัก โดยให้น้ำหนักภาพร้อยละ 60 และข้อความร้อยละ 40 ซึ่งพิสูจน์แล้วว่าช่วยให้ระบบทนทานต่อข้อมูลอาการที่ขัดแย้งกับภาพได้อย่างดีเยี่ยม
\end{enumerate}
% ------------------------------------------------------------------------------
% 2.6 การทบทวนงานวิจัยและปริญญานิพนธ์ที่เกี่ยวข้อง
% ------------------------------------------------------------------------------
\section{การทบทวนงานวิจัยและปริญญานิพนธ์ที่เกี่ยวข้อง}
\label{sec:literature_review_works}

\subsection{ปริญญานิพนธ์ของรุ่นพี่กลุ่มที่ 1 (ปีการศึกษา 2567)}
\textbf{เปี่ยมเดช ยศเงินทองฟู และ ปุณณวิช สิทธิเสือ (2567) \cite{piamdech2024oral}} ได้พัฒนาปริญญานิพนธ์เรื่อง \textit{``การพัฒนาระบบเว็บผู้เชี่ยวชาญและส่วนต่อประสานแอปพลิเคชันไลน์เพื่อการประเมินความเสี่ยงมะเร็งในช่องปากจากรูปแผลและประวัติผู้ป่วย''} โดยใช้แบบจำลอง YOLO (YOLOv8, YOLOv9, YOLOv10) เพื่อตรวจจับและจำแนกรอยโรคในช่องปาก 5 ชนิด จากชุดข้อมูลภาพถ่ายของคณะทันตแพทยศาสตร์ มหาวิทยาลัยมหิดล จำนวน 1,579 ภาพ ควบคู่กับการนำประวัติผู้ป่วยมาประเมินร่วมผ่านกฎเงื่อนไข (Rule-based) บนระบบคลาวด์ GCP และ LINE API
\begin{itemize}
    \item \textbf{ข้อค้นพบและจุดเด่น:} แบบจำลอง YOLOv8 ให้ความถูกต้องสูงสุดที่ร้อยละ 86 เมื่อนำข้อมูลประวัติผู้ป่วยมาประเมินร่วม และมีจุดเด่นคือสามารถวาดกรอบ Bounding Box ระบุตำแหน่งรอยโรคบนภาพได้ชัดเจน
    \item \textbf{ข้อจำกัดและปัญหาที่พบ:} การทำงานระหว่างการตรวจจับภาพและกฎเงื่อนไขอาการเป็นการประมวลผลแบบแยกส่วน (Decoupled) หากภาพถ่ายไม่ชัดจนโมเดลเกิดภาวะ No Detection ระบบจะไม่สามารถใช้อาการของผู้ป่วยมาช่วยวิเคราะห์ต่อได้เลย นอกจากนี้ การตัดสินใจด้วยกฎเงื่อนไขยังไม่สามารถรับมือกับข้อมูลอาการที่มีความขัดแย้งกันได้ดีพอ
\end{itemize}

\subsection{ปริญญานิพนธ์ของรุ่นพี่กลุ่มที่ 2 (ปีการศึกษา 2568)}
\textbf{รตา อุณหะ และ วรดา สาระคนธ์ (2568) \cite{unnahasarakhon2025multimodal}} ได้พัฒนาปริญญานิพนธ์เรื่อง \textit{``การพัฒนาอัลกอริทึมสำหรับระบบผู้เชี่ยวชาญวิเคราะห์รอยโรคทางช่องปากผ่าน Web และ Line API''} โดยมุ่งเน้นการแก้ปัญหาการประมวลผลแบบแยกส่วนของรุ่นพี่กลุ่มที่ 1 ด้วยการพัฒนาแบบจำลองการเรียนรู้เชิงลึกแบบหลายโมดาลิตี (Multimodal Deep Learning) ที่ผสานภาพถ่าย (DenseNet-121) เข้ากับข้อความประวัติผู้ป่วยภาษาไทย (WangchanBERTa) และเปรียบเทียบวิธีการรวมข้อมูลทั้ง Early Fusion และ Late Fusion
\begin{itemize}
    \item \textbf{ข้อค้นพบและจุดเด่น:} วิธีการ Late Fusion แบบกำหนดน้ำหนักภาพต่อข้อความ 60:40 ให้ประสิทธิภาพสูงสุด โดยมีค่าความแม่นยำถึงร้อยละ 95.5 และยังคงรักษาระดับความแม่นยำได้ที่ร้อยละ 82.5 ในสถานการณ์จำลองที่ข้อมูลอาการไม่ตรงกับภาพ นอกจากนี้ยังมีภาพความร้อน Grad-CAM อธิบายการตัดสินใจ
    \item \textbf{ข้อจำกัดและปัญหาที่พบ:} แบบจำลองเป็นแบบการจำแนกภาพทั้งภาพ (Whole-image classification) ทำให้ไม่สามารถตีกรอบ Bounding Box ชี้ตำแหน่งรอยโรคได้เหมือนระบบ YOLO นอกจากนี้ ด้านสถาปัตยกรรมระบบยังรันแบบต้นแบบ ทำให้กินทรัพยากรสูง และยังขาดระบบแดชบอร์ดสำหรับให้แพทย์ส่งข้อมูลย้อนกลับ (Feedback Loop)
\end{itemize}

\subsection{งานวิจัยการตรวจคัดกรองมะเร็งช่องปากแบบหลายโมดาลิตีในระดับสากล}
\textbf{Devindi et al. (IEEE Access, 2024, เอกสารเลขที่ 10664507) \cite{devindi2024multimodal}} ได้นำเสนอแนวคิด \textit{``Multimodal Deep Convolutional Neural Network Pipeline for AI-Assisted Early Detection of Oral Cancer''} ซึ่งสอดคล้องกับแนวคิดของโครงงานนี้ โดยมุ่งเน้นการบูรณาการข้อมูลภาพถ่ายจากสมาร์ตโฟนเข้ากับข้อมูลอภิพันธุ์ของผู้ป่วย (Patient Metadata เช่น ข้อมูลประวัติและอาการทางคลินิก) เพื่อเลียนแบบกระบวนการตรวจวินิจฉัยจริงของแพทย์ งานวิจัยได้เปรียบเทียบสถาปัตยกรรมตัวเข้ารหัสภาพ 6 แบบ (MobileNetV3-Large, MixNet-S, ResNet-50, HRNet-W18-C, DenseNet-121 และ Inception\_v3) ผลการวิจัยยืนยันว่าการใช้ข้อมูลภาพถ่ายร่วมกับประวัติผู้ป่วยให้ค่า Accuracy ร้อยละ 81 และ F1-score ร้อยละ 78 ซึ่งเหนือกว่าการใช้ภาพถ่ายเพียงอย่างเดียวอย่างมีนัยสำคัญ งานวิจัยนี้เป็นหลักฐานเชิงประจักษ์ระดับนานาชาติที่ยืนยันว่าการรวมข้อมูลภาพและประวัติคือทิศทางที่ถูกต้องในการคัดกรองมะเร็งช่องปาก

\subsection{การทบทวนวรรณกรรมอย่างเป็นระบบด้านการเรียนรู้เชิงลึกในมะเร็งช่องปาก}
\textbf{Warin และ Suebnukarn (BMC Oral Health, 2024, DOI: 10.1186/s12903-024-03993-5) \cite{warin2024deep}} ได้ทำการศึกษาทบทวนวรรณกรรมอย่างเป็นระบบตามแนวทางมาตรฐาน PRISMA ครอบคลุมงานวิจัยด้าน Deep Learning ในมะเร็งช่องปากทั่วโลก ผลการศึกษาชี้ให้เห็นว่าสถาปัตยกรรม CNN และ YOLO มีประสิทธิภาพสูงในการคัดกรอง แต่จุดอ่อนสำคัญของงานวิจัยในปัจจุบันคือการขาดการนำไปใช้งานจริงในคลินิก (Clinical Translation) เนื่องจากขาดระบบเชื่อมต่อที่ง่าย ขาดความสามารถในการอธิบายผล และที่สำคัญคือขาดการผสานข้อมูลบริบททางคลินิกของผู้ป่วย ซึ่งข้อเสนอแนะของงานวิจัยนี้ตรงกับวัตถุประสงค์ในการสร้างแพลตฟอร์มสนับสนุนการตัดสินใจของโครงงานเรา

\begin{table}[H]
\centering
\small
\caption{ตารางสังเคราะห์เปรียบเทียบวรรณกรรมและปริญญานิพนธ์ที่เกี่ยวข้อง}
\label{tab:synthesis_matrix}
\begin{tabular}{p{2.8cm}p{2.2cm}p{2.8cm}p{2.2cm}p{3.5cm}}
\toprule
\textbf{งานวิจัย/ปริญญานิพนธ์} & \textbf{ข้อมูลนำเข้า} & \textbf{สถาปัตยกรรมโมเดล} & \textbf{รอยโรคเป้าหมาย} & \textbf{ผลลัพธ์และข้อสังเกต} \\
\midrule
เปี่ยมเดช และ ปุณณวิช (2567) \cite{piamdech2024oral} & ภาพถ่าย + กฎเงื่อนไข & YOLOv8/v9/v10 + Rule-based & 5 รอยโรค & ตีกรอบ Bounding Box ได้ แต่ระบบแยกส่วนและเสี่ยง No Detection \\
รตา และ วรดา (2568) \cite{unnahasarakhon2025multimodal} & ภาพ + ประวัติไทย & DenseNet-121 + WangchanBERTa & 5 รอยโรค & แม่นยำ 95.5\% (Late Fusion) แต่ไม่มี Bounding Box และเป็นงานต้นแบบ \\
Devindi et al. (2024) \cite{devindi2024multimodal} & ภาพ + Metadata & MobileNetV3, DenseNet, Grad-CAM & Benign / Malignant & Accuracy 81\%, ยืนยันว่า Multimodal เหนือกว่าภาพเดี่ยว \\
Warin \& Suebnukarn (2024) \cite{warin2024deep} & การทบทวนวรรณกรรม & Systematic Review (PRISMA) & OSCC, OPMD & ชี้ชัดถึงความจำเป็นของ Multimodal, ระบบจริง และ Explainable AI \\
\textbf{โครงงานนี้ (2569)} & \textbf{ภาพ + ประวัติไทย + Feedback แพทย์} & \textbf{Hybrid: YOLO + DenseNet-BERT} & \textbf{5 รอยโรค + มะเร็งช่องปาก} & \textbf{Production Platform: มีทั้ง Bounding Box, Multimodal Fusion, Doctor Dashboard, CI/CD, PDPA} \\
\bottomrule
\end{tabular}
\end{table}

% ------------------------------------------------------------------------------
% 2.7 สรุปแนวทางการออกแบบสถาปัตยกรรมและกรอบแนวคิดของโครงงาน
% ------------------------------------------------------------------------------
\section{สรุปแนวทางการออกแบบสถาปัตยกรรมและกรอบแนวคิดของโครงงาน}
\label{sec:system_framework_summary}

จากการทบทวนวรรณกรรมและข้อจำกัดของโครงงานรุ่นพี่ทั้ง 2 รุ่น คณะผู้จัดทำจึงได้กำหนดกรอบแนวคิดในการพัฒนาระบบของโครงงานนี้ ให้ก้าวข้ามจากระดับงานวิจัยต้นแบบไปสู่ แพลตฟอร์มสนับสนุนการวินิจฉัยรอยโรคทางช่องปากที่พร้อมใช้งานจริงสำหรับทันตแพทย์ (Production-Ready Clinical Decision Support Platform) โดยมีหลักการสำคัญดังนี้:

\begin{enumerate}
    \item \textbf{การบูรณาการแบบจำลองปัญญาประดิษฐ์แบบทางเลือก (Alternative AI Model Integration):} นำแบบจำลอง YOLO สำหรับการตรวจจับและระบุตำแหน่งรอยโรค (Bounding Box) และแบบจำลองแบบหลายโหมด (Multimodal Model) มาประยุกต์ใช้งานร่วมกันภายใต้สถาปัตยกรรมไมโครเซอร์วิส (Microservices Architecture) เพื่อเพิ่มความยืดหยุ่นให้ผู้ใช้งานสามารถเลือกใช้แบบจำลองได้ตามความเหมาะสม และปรับปรุง User Interface ให้ทันสมัยและใช้งานง่ายมากขึ้น
    \item \textbf{การผสมผสานสองระบบปัญญาประดิษฐ์ (Hybrid Inference Engine):} บูรณาการทั้งโมเดล YOLO สำหรับการระบุตำแหน่งรอยโรคด้วย Bounding Box และโมเดล Multimodal สำหรับการประเมินคะแนนความเสี่ยงมะเร็งที่แม่นยำสูงเข้าด้วยกัน
    
    \item \textbf{การสร้างระบบวงจรข้อมูลย้อนกลับ (Doctor Feedback Loop):} สร้างแดชบอร์ดให้ทันตแพทย์สามารถตรวจสอบและแก้ไขผลการทำนายของ AI ได้ เพื่อเปลี่ยนแพลตฟอร์มให้กลายเป็นแหล่งรวบรวมข้อมูล Ground Truth ทางคลินิกคุณภาพสูงสำหรับการพัฒนาโมเดลต่อไป
    \item \textbf{การออกแบบสถาปัตยกรรมไมโครเซอร์วิสและ CI/CD:} แยกส่วน API Gateway (FastAPI) และ AI Engine ผ่าน Docker Containers รองรับการขยายตัวอัตโนมัติบนคลาวด์ พร้อมทั้งมีระบบ Automated Testing และ Continuous Deployment เพื่อความเสถียรสูงสุด
    \item \textbf{การรักษาความปลอดภัยและคุ้มครองข้อมูลส่วนบุคคล (PDPA Compliance):} วางระบบยืนยันตัวตน การกำหนดสิทธิ์ตามบทบาท (RBAC) และการเข้ารหัสข้อมูล เพื่อให้การประมวลผลข้อมูลสุขภาพของผู้ป่วยเป็นไปตามกฎหมายอย่างเคร่งครัด
\end{enumerate}



    % ==============================================================================
% บทที่ 3: วิธีการดำเนินโครงงาน (Research Methodology and System Implementation)
% โครงการ: การพัฒนาแพลตฟอร์มสำหรับการสนับสนุนการวินิจฉัยรอยโรคทางช่องปาก
%          และผสมผสานการใช้งานระบบผู้เชี่ยวชาญสำหรับทันตแพทย์
% ==============================================================================

\chapter{วิธีการดำเนินโครงงาน}
\label{chapter:methodology}

โครงงานนี้มีเป้าหมายในการพัฒนาแพลตฟอร์มเว็บแอปพลิเคชันเพื่อสนับสนุนการตัดสินใจทางคลินิกสำหรับทันตแพทย์ (Clinical Decision Support Platform) โดยมุ่งเน้นการนำแบบจำลองปัญญาประดิษฐ์ที่ได้รับการวิจัยและพัฒนาขึ้นจากโครงงานของรุ่นพี่ในปีก่อนหน้าทั้ง 2 กลุ่ม ได้แก่ แบบจำลองตรวจจับวัตถุ (YOLO) และแบบจำลองการเรียนรู้เชิงลึกแบบหลายโมดาลิตี (Multimodal Deep Learning: DenseNet-121 ร่วมกับ WangchanBERTa) นำมารวบรวมและบูรณาการเข้าด้วยกันภายใต้สถาปัตยกรรมระบบไมโครเซอร์วิส เพื่อให้ทันตแพทย์และบุคลากรทางการแพทย์สามารถเลือกใช้งานแบบจำลอง (Alternative Model Selection) ให้เหมาะสมกับบริบทและความต้องการในการตรวจคัดกรองรอยโรคแต่ละกรณี พร้อมทั้งพัฒนาส่วนติดต่อผู้ใช้งานที่ทันสมัยและสะดวกต่อการปฏิบัติงานจริง

ขั้นตอนและระเบียบวิธีวิจัยในการดำเนินโครงงานแบ่งออกเป็น 6 ส่วนหลัก ดังนี้:
\begin{enumerate}
    \item การรวบรวมและวิเคราะห์คุณลักษณะของแบบจำลองปัญญาประดิษฐ์จากงานวิจัยรุ่นพี่
    \item การออกแบบสถาปัตยกรรมระบบแบบไมโครเซอร์วิส 3 คอนเทนเนอร์ (System Architecture \& Microservices)
    \item การออกแบบผังกรณีการใช้งานและผังกิจกรรมการทำงานของระบบ (Use Case \& Activity Diagram)
    \item การพัฒนาส่วนติดต่อผู้ใช้งานสำหรับทันตแพทย์ (User Interface Design)
    \item การจัดเตรียมสภาพแวดล้อม การคอนเทนเนอร์ไรซ์ และการทดสอบระบบอัตโนมัติ
    \item สรุปกรอบการดำเนินงานเพื่อเตรียมการประเมินผลในบทที่ 4
\end{enumerate}

% ------------------------------------------------------------------------------
% 3.1 การรวบรวมและวิเคราะห์แบบจำลองปัญญาประดิษฐ์จากงานวิจัยรุ่นพี่
% ------------------------------------------------------------------------------
\section{การรวบรวมและวิเคราะห์แบบจำลองปัญญาประดิษฐ์จากงานวิจัยรุ่นพี่}
\label{sec:ai_model_consolidation}

ในการดำเนินงานขั้นแรก คณะผู้จัดทำได้นำไฟล์โมเดลที่ผ่านการฝึกฝน (Pre-trained Weights), ซอร์สโค้ดการอนุมานผล (Inference Pipeline), และชุดกฎเกณฑ์ทางคลินิกจากปริญญานิพนธ์ของรุ่นพี่ทั้งสองกลุ่ม มารวบรวมและวิเคราะห์โครงสร้างข้อมูลนำเข้าและส่งออก (Input/Output Interface) เพื่อเตรียมการผสานรวมเข้าสู่ระบบ:

\subsection{แบบจำลองตรวจจับวัตถุ YOLO จากรุ่นพี่กลุ่มที่ 1 (ปีการศึกษา 2567)}
แบบจำลองของรุ่นพี่กลุ่มที่ 1 (เปี่ยมเดช และ ปุณณวิช) \cite{piamdech2024oral} เป็นแบบจำลองในกลุ่ม Object Detection ได้แก่ สถาปัตยกรรม YOLOv7 (\texttt{yolov7.pt}) และ YOLOv8 (\texttt{yolov8.pt}) ซึ่งมีลักษณะการทำงานสำคัญดังนี้:
\begin{enumerate}
    \item \textbf{ข้อมูลนำเข้า (Input):} ภาพถ่ายรอยโรคในช่องปากที่ผ่านการแปลงเป็นปริภูมิสี BGR และปรับขนาดภาพด้วยวิธี Letterbox ขนาด $640 \times 640$ พิกเซล ควบคู่กับข้อความอาการของผู้ป่วย
    \item \textbf{การตรวจจับและคัดกรองกรอบรอยโรค:} แบบจำลองจะทำการสกัดคุณลักษณะและทำนายพิกัด Bounding Box ในรูปแบบ $[x, y, w, h]$ พร้อมคะแนนความเชื่อมั่น โดยตั้งเกณฑ์คัดกรองขั้นต่ำ (Confidence Threshold) ไว้ที่ 0.25 จากนั้นใช้กระบวนการคัดกรองกรอบรอยโรคที่ซ้อนทับกัน (Concentric Box \& Overlap Suppression) โดยหากกรอบของคลาสเดียวกันมีค่าความทับซ้อน (IoU) เกินร้อยละ 40 หรือมีพื้นที่ครอบคลุมอยู่ในอีกกรอบหนึ่งเกินร้อยละ 70 จะทำการตัดกรอบที่มีคะแนนน้อยกว่าออกไป
    \item \textbf{การจับคู่กฎเกณฑ์อาการทางคลินิก (Clinical Rule Matching):} นำข้อความอาการของผู้ป่วยมาตัดแยกคำสำคัญ (Tokens) เพื่อเปรียบเทียบกับพจนานุกรมอาการเฉพาะของรอยโรค (\texttt{PATIENT\_CONDITIONS}) เช่น หากตรวจพบรอยโรค Leukoplakia และพบคำสำคัญ ``สูบบุหรี่'', ``ดื่มเหล้า'', หรือ ``เคี้ยวหมาก'' จะถือว่าอาการตรงกัน
    \item \textbf{การคำนวณคะแนนความมั่นใจขั้นสุดท้าย:}
    \begin{itemize}
        \item กรณีอาการตรงกับรอยโรค: ปรับคะแนนเพิ่มตามสูตร $\text{Score} = \frac{\text{conf}}{2} + 0.50$
        \item กรณีอาการไม่ตรงกับรอยโรค: ปรับคะแนนลดลงตามสูตร $\text{Score} = \frac{\text{conf}}{2}$ (หากเดิมเกิน 0.60) หรือคงค่าเดิมไว้
    \end{itemize}
    \item \textbf{ผลลัพธ์ที่ได้ (Output):} ภาพถ่ายที่ผ่านการวาดกรอบ Bounding Box พร้อมป้ายกำกับชื่อโรค ค่าเปอร์เซ็นต์ความเชื่อมั่น และสีเส้นกรอบจำเพาะตามชนิดรอยโรค
\end{enumerate}

\subsection{แบบจำลองการเรียนรู้เชิงลึกแบบหลายโมดาลิตีจากรุ่นพี่กลุ่มที่ 2 (ปีการศึกษา 2568)}
แบบจำลองของรุ่นพี่กลุ่มที่ 2 (รตา และ วรดา) \cite{unnahasarakhon2025multimodal} พัฒนาขึ้นบนแนวคิด Multimodal Deep Learning ซึ่งผสานแบบจำลองภาพและข้อความเข้าด้วยกัน ประกอบด้วย:
\begin{enumerate}
    \item \textbf{โมดูลประมวลผลภาพ (DenseNet-121):} โหลดน้ำหนักที่ผ่านการเทรน (\texttt{best\_densenet121\_img.pth}) รับภาพขนาด $600 \times 600$ พิกเซล และทำ Normalization ด้วยค่าเฉลี่ยและส่วนเบี่ยงเบนมาตรฐานจำเพาะของชุดข้อมูลรอยโรคช่องปาก สกัดเวกเตอร์คุณลักษณะขนาด $1 \times 1024$ มิติ ผ่านชั้น Global Average Pooling (GAP)
    \item \textbf{โมดูลประมวลผลภาษาไทย (WangchanBERTa):} โหลดน้ำหนัก (\texttt{best\_text.pth}) รับข้อความอาการภาษาไทย เข้าสู่ SentencePiece Tokenizer เพื่อแปลงเป็น Input IDs และ Attention Mask แล้วส่งผ่าน Transformer Encoder เพื่อสกัดเวกเตอร์ตัวแทนบริบทจากโทเคน \texttt{[CLS]} ขนาด $1 \times 768$ มิติ
    \item \textbf{การรวมผลลัพธ์ระดับการตัดสินใจ (Late Decision Fusion):} นำคะแนนความน่าจะเป็นของโมเดลภาพ ($P_{img}$) และโมเดลข้อความ ($P_{text}$) มารวมกันด้วยการถ่วงน้ำหนักภาพร้อยละ 60 และข้อความร้อยละ 40:
    \begin{equation}
        P_{final} = 0.60 \times P_{img} + 0.40 \times P_{text}
        \label{eq:ch3_late_fusion}
    \end{equation}
    \item \textbf{การสร้างภาพความร้อนอธิบายผล (Grad-CAM Heatmap):} คำนวณเกรเดียนต์จากเลเยอร์สุดท้ายของ DenseNet-121 เพื่อสร้างแผนที่ความร้อนแสดงตำแหน่งรอยโรคที่ AI ใช้ในการตัดสินใจ และนำไปทับซ้อนลงบนภาพจริง
    \item \textbf{ผลลัพธ์ที่ได้ (Output):} ชื่อรอยโรคที่ทำนาย, เปอร์เซ็นต์ความเชื่อมั่นรวม, และภาพผลลัพธ์ที่ซ้อนทับด้วย Grad-CAM Heatmap
\end{enumerate}

\subsection{การเปรียบเทียบคุณลักษณะของแบบจำลองทั้งสองระบบ}
เพื่อให้ผู้ใช้งานเข้าใจความแตกต่างและสามารถเลือกใช้งานได้อย่างถูกต้อง ตารางที่ \ref{tab:model_comparison} แสดงการเปรียบเทียบจุดเด่นและคุณลักษณะเชิงเทคนิคของทั้งสองแบบจำลอง:

\begin{table}[htbp]
\centering
\small
\caption{การเปรียบเทียบคุณสมบัติเชิงเทคนิคระหว่างแบบจำลอง YOLO และแบบจำลอง Multimodal}
\label{tab:model_comparison}
\begin{tabular}{p{3.5cm}p{5.0cm}p{5.0cm}}
\toprule
\textbf{มิติการเปรียบเทียบ} & \textbf{แบบจำลอง YOLO (รุ่นพี่กลุ่ม 1)} & \textbf{แบบจำลอง Multimodal (รุ่นพี่กลุ่ม 2)} \\
\midrule
สถาปัตยกรรมหลัก & YOLOv7 / YOLOv8 & DenseNet-121 + WangchanBERTa \\
รูปแบบการทำงาน & การตรวจจับวัตถุ (Object Detection) & การจำแนกภาพแบบหลายโหมด (Multimodal) \\
การระบุตำแหน่งรอยโรค & \textbf{ระบุได้ด้วย Bounding Box ชัดเจน} & ระบุด้วยภาพความร้อน Grad-CAM Heatmap \\
การผสานข้อมูลประวัติ & ใช้กฎเงื่อนไขภายนอก (Rule-based) & \textbf{รวมด้วย Late Fusion ภายในโมเดล (60:40)} \\
ความแม่นยำทางคลินิก & ร้อยละ 80--86 & \textbf{ร้อยละ 95.5 (ทนทานต่อข้อมูลขัดแย้ง)} \\
การจัดการกรณีภาพไม่ชัด & เสี่ยงต่อภาวะ No Detection & ยังคงสามารถใช้ข้อมูลข้อความช่วยประเมินได้ \\
ความเหมาะสมในการใช้งาน & \textbf{เหมาะสำหรับการตรวจหาตำแหน่งแผล} & \textbf{เหมาะสำหรับการประเมินความเสี่ยงมะเร็ง} \\
\bottomrule
\end{tabular}
\end{table}

\subsection{ยุทธศาสตร์และเหตุผลในการรวมแบบจำลองเข้าด้วยกัน (Model Consolidation Strategy)}
จากการวิเคราะห์แบบจำลองทั้งสองกลุ่ม คณะผู้จัดทำพบว่าไม่มีแบบจำลองใดสมบูรณ์แบบในทุกมิติ โดยแต่ละแบบจำลองมีข้อได้เปรียบที่ส่งเสริมซึ่งกันและกัน (Complementary Strengths):
\begin{enumerate}
    \item \textbf{การเสริมจุดเด่นด้านตำแหน่งและบริบท (Localization vs Contextual Understanding):} แบบจำลอง YOLO มีจุดแข็งในการระบุพิกัดและขนาดของรอยโรคอย่างชัดเจนผ่าน Bounding Box ซึ่งจำเป็นต่อการวางแผนตัดชิ้นเนื้อ (Biopsy) ในขณะที่แบบจำลอง Multimodal มีความโดดเด่นในการประมวลผลข้อความอาการและพฤติกรรมเสี่ยงของผู้ป่วย เช่น การสูบบุหรี่จัด การดื่มสุรา หรืออาการเจ็บแสบ ทำให้ประเมินความเสี่ยงมะเร็งได้รอบด้านกว่า
    \item \textbf{ความยืดหยุ่นในการสนับสนุนการตัดสินใจทางคลินิก (Clinical Workflow Flexibility):} ในบริบทคลินิกจริง ผู้ป่วยแต่ละรายมีอาการและลักษณะรอยโรคที่แตกต่างกัน หากรอยโรคมีขอบเขตชัดเจน ทันตแพทย์อาจต้องการดูขอบเขตแผลด้วยโมเดล YOLO แต่หากรอยโรคมีลักษณะก้ำกึ่ง แบนราบ หรือมีความเสี่ยงแฝง ทันตแพทย์จะต้องการผลการประเมินจาก Multimodal เพื่อดูภาพความร้อน Grad-CAM และระดับความเสี่ยง
    \item \textbf{การตรวจสอบยืนยันข้ามแบบจำลองเพื่อความเห็นที่สอง (Cross-Model Validation / Second Opinion):} แพลตฟอร์มที่รวมทั้งสองแบบจำลองไว้ในระบบเดียวกัน เปิดโอกาสให้ทันตแพทย์สามารถส่งตรวจภาพและอาการเคสเดียวกันผ่านทั้งสองแบบจำลองเพื่อเปรียบเทียบผลลัพธ์ หากผลการวิเคราะห์สอดคล้องกัน จะช่วยเพิ่มความมั่นใจในการวินิจฉัยอย่างมีนัยสำคัญ
\end{enumerate}

ดังนั้น เป้าหมายหลักของโครงงานนี้จึงไม่ใช่การฝึกแบบจำลองใหม่ตั้งแต่ต้น แต่เป็นการสร้าง \textbf{Consolidated AI Platform} ที่นำแบบจำลองของรุ่นพี่ทั้งสองชุดมารวมเป็นระบบเดียว ออกแบบกลไกการสลับและเลือกใช้งานแบบจำลอง (Alternative Model Selection) และแปลงผลลัพธ์ให้อยู่ในมาตรฐานเดียวกัน (Unified Schema) เพื่อพร้อมสำหรับการใช้งานจริงในคลินิกทันตกรรม

% ------------------------------------------------------------------------------
% 3.2 การออกแบบสถาปัตยกรรมระบบและการบูรณาการแบบจำลอง
% ------------------------------------------------------------------------------
\section{การออกแบบสถาปัตยกรรมระบบและการบูรณาการแบบจำลอง}
\label{sec:system_architecture_design}

\subsection{สถาปัตยกรรมไมโครเซอร์วิส (Microservices Architecture)}
เพื่อแก้ไขปัญหาการผูกติดกันของระบบและลดการกินทรัพยากรของสถาปัตยกรรมแบบ Monolithic เดิม ระบบใหม่ได้รับการออกแบบตามสถาปัตยกรรม Microservices โดยแบ่งการทำงานออกเป็น 3 คอนเทนเนอร์หลัก (Three-Tier Microservices) ที่แยกหน้าที่กันอย่างอิสระและสื่อสารกันผ่านระบบเครือข่ายภายใน (\texttt{oralscan-network}):
\begin{enumerate}
    \item \textbf{Frontend Web Service (\texttt{oralscan-frontend}, พอร์ต 5173):} ส่วนติดต่อผู้ใช้งานฝั่งหน้าบ้าน พัฒนาด้วย React ร่วมกับ Vite และ Tailwind CSS ออกแบบในรูปแบบ Single Page Application (SPA) เพื่อให้ทันตแพทย์สามารถใช้งานได้อย่างสะดวก รวดเร็ว รองรับการอัปโหลดรูปภาพแบบ Drag-and-Drop มีระบบเลือกโมเดล (Model Selector) และเลือกอาการทางคลินิก (Symptoms Checkbox) โดยสื่อสารกับหลังบ้านผ่าน RESTful API
    \item \textbf{Backend API Gateway Service (\texttt{oralscan-backend}, พอร์ต 8000):} ส่วนหลังบ้านที่พัฒนาด้วย FastAPI ทำหน้าที่เป็นจุดรับคำขอภายนอกเพียงจุดเดียว (Single Entry Point) คอยตรวจสอบความถูกต้องของข้อมูลนำเข้า (Data Validation) ผ่าน Pydantic Schema, จัดการระบบรักษาความปลอดภัยและการจำกัดสิทธิ์ (CORS), จัดเก็บและให้บริการไฟล์รูปภาพทั้งภาพต้นฉบับและภาพผลลัพธ์ใน Static Directory (\texttt{/app/static}), จัดการเส้นทางคำขอ (Dynamic Model Router), และแปลงผลลัพธ์จาก AI ให้อยู่ในมาตรฐานเดียวกัน (Unified Response) พร้อมประเมินระดับความเสี่ยงมะเร็งทางคลินิก
    \item \textbf{AI Inference Engine Service (\texttt{oralscan-ai-engine}, พอร์ต 8001):} เซอร์วิสประมวลผลโมเดลปัญญาประดิษฐ์ความเร็วสูง บรรจุสภาพแวดล้อมทางคณิตศาสตร์และแมชชีนเลิร์นนิง ได้แก่ PyTorch, Ultralytics YOLO, OpenCV-headless และ HuggingFace Transformers โดยใช้ Thread-Safe Singleton Pattern ในการโหลดน้ำหนักของโมเดลรุ่นพี่ทั้งสองกลุ่มขึ้นสู่หน่วยความจำ (ได้แก่ \texttt{YoloService} และ \texttt{MultimodalService}) เพื่อป้องกันปัญหาโหลดโมเดลซ้ำซ้อน ประหยัด RAM และรองรับการประมวลผลคำขอพร้อมกันอย่างมีเสถียรภาพ
\end{enumerate}

\subsection{รูปแบบการเลือกใช้งานแบบจำลอง (Alternative Model Selection)}
เมื่อมีคำขอส่งตรวจเข้ามาผ่านเอนด์พอยต์ \texttt{/api/assess} ผู้ใช้งานสามารถส่งค่าพารามิเตอร์ \texttt{modelType} เพื่อกำหนดเส้นทางการประมวลผล (Dynamic Model Routing) ได้ตามความต้องการทางคลินิก ดังแสดงในรูปที่ \ref{fig:model_pipeline_diagram}:
\begin{itemize}
    \item \textbf{กรณีเลือกโมเดล YOLO (หรือ YOLOv7 / YOLOv8):} API Gateway จะส่งงานต่อไปยังฟังก์ชัน \texttt{infer\_yolo} ภายใน AI Engine โดยนำภาพรอยโรคเข้าสู่โมเดลตรวจจับวัตถุเพื่อค้นหา Bounding Box พร้อมทั้งนำอาการที่ผู้ใช้กรอกเข้าสู่กลไกจับคู่กฎเกณฑ์ทางคลินิก (Clinical Rule Matching) เพื่อคำนวณคะแนนความเชื่อมั่นสุทธิ และวาดกรอบสีจำเพาะลงบนภาพ
    \item \textbf{กรณีเลือกโมเดล DenseNet-BERT (Multimodal):} API Gateway จะส่งงานไปยังฟังก์ชัน \texttt{infer\_densenet\_bert} ภายใน AI Engine โดยแยกการประมวลผลภาพผ่าน DenseNet-121 และข้อความอาการผ่าน WangchanBERTa จากนั้นรวมผลด้วยกลไก Late Decision Fusion (น้ำหนักภาพร้อยละ 60 และข้อความร้อยละ 40) พร้อมสร้างภาพความร้อน Grad-CAM Heatmap
\end{itemize}

% แผนภาพขั้นตอนการรวมและการเลือกแบบจำลอง AI (เว้นกรอบ Placeholder สำหรับใส่ภาพกราฟิก)
\begin{figure}[htbp]
\centering
\fbox{
    \parbox{0.92\textwidth}{
        \centering
        \vspace{2.5cm}
        \textbf{[พื้นที่สำหรับใส่ภาพ: แผนภาพขั้นตอนการบูรณาการและการเลือกแบบจำลอง AI (AI Model Consolidation \& Selection Pipeline)]}\\[0.3cm]
        {\footnotesize \textit{(Inference Pipeline: Image \& Symptoms $\rightarrow$ Model Selection Router $\rightarrow$ [Branch A: YOLOv7/v8 + Rule Matching] vs [Branch B: Multimodal DenseNet-121 + WangchanBERTa] $\rightarrow$ Unified Schema Output)}}\\[0.2cm]
        {\scriptsize สามารถใส่ภาพ \texttt{images/model\_selection\_pipeline.png} ลงในจุดนี้}
        \vspace{2.5cm}
    }
}
\caption{แผนภาพขั้นตอนการบูรณาการและการเลือกแบบจำลองปัญญาประดิษฐ์ (AI Model Consolidation and Selection Pipeline)}
\label{fig:model_pipeline_diagram}
\end{figure}

\subsection{การกำหนดรูปแบบข้อมูลผลลัพธ์มาตรฐาน (Unified Assessment Schema)}
ไม่ว่าผู้ใช้งานจะเลือกประมวลผลด้วยแบบจำลองใด API Gateway จะส่งคืนผลลัพธ์ในโครงสร้าง JSON มาตรฐานเดียวกันเสมอ (\texttt{UnifiedAssessmentResponse}) ซึ่งประกอบด้วย:
\begin{itemize}
    \item \texttt{model\_used}: สถาปัตยกรรมโมเดลที่ใช้งานจริงในการตรวจ
    \item \texttt{disease\_name}: ชื่อรอยโรคเป้าหมายที่ตรวจพบ (Ulcer, Candidiasis, Linea alba buccalis, Leukoplakia, หรือ Lichen planus)
    \item \texttt{risk\_level}: ระดับความเสี่ยงต่อมะเร็งช่องปาก (Low, Medium, High)
    \item \texttt{confidence\_score}: ค่าคะแนนความเชื่อมั่นของแบบจำลอง (ทศนิยม 0.00 ถึง 1.00)
    \item \texttt{bounding\_boxes}: พิกัดกรอบรอยโรค (มีข้อมูลเมื่อใช้โมเดล YOLO และเป็นค่าว่างเมื่อใช้ DenseNet-BERT)
    \item \texttt{processed\_image\_url}: URL สำหรับดึงภาพที่ผ่านการวาดกรอบหรือ Grad-CAM ไปแสดงผล
\end{itemize}

\subsection{เกณฑ์การประเมินระดับความเสี่ยงมะเร็งทางคลินิก (Clinical Risk Level)}
เพื่อช่วยให้ทันตแพทย์จัดลำดับความเร่งด่วนในการรักษาได้ทันที ระบบได้กำหนดฟังก์ชันประเมินระดับความเสี่ยงมะเร็ง $\text{RiskLevel}(D, C)$ โดยพิจารณาจากชนิดรอยโรค ($D$) และคะแนนความเชื่อมั่น ($C$):
\begin{equation}
    \text{RiskLevel}(D, C) = 
    \begin{cases}
        \text{High (สูง)} & \text{เมื่อ } D \in \{\text{Leukoplakia}, \text{Ulcer}\} \text{ และ } C \ge 0.40 \\
        \text{Medium (ปานกลาง)} & \text{เมื่อ } D \in \{\text{Leukoplakia}, \text{Ulcer}\} \text{ และ } C < 0.40 \\
        \text{Medium (ปานกลาง)} & \text{เมื่อ } D \in \{\text{Lichen Planus}\} \\
        \text{Low (ต่ำ)} & \text{เมื่อ } D \in \{\text{Candidiasis}, \text{Linea Alba}\} \\
        \text{Low (ต่ำ)} & \text{เมื่อ } D \in \{\text{No Detection}, \text{Normal}\}
    \end{cases}
    \label{eq:ch3_risk_level}
\end{equation}

% แผนภาพสถาปัตยกรรมระบบ (เว้นกรอบ Placeholder สำหรับใส่ภาพกราฟิก)
\begin{figure}[htbp]
\centering
\fbox{
    \parbox{0.92\textwidth}{
        \centering
        \vspace{2.5cm}
        \textbf{[พื้นที่สำหรับใส่ภาพ: แผนภาพสถาปัตยกรรมระบบไมโครเซอร์วิส 3 คอนเทนเนอร์]}\\[0.3cm]
        {\footnotesize \textit{(System Architecture Diagram: Frontend Service (Port 5173) $\leftrightarrow$ Backend API Gateway (Port 8000) $\leftrightarrow$ AI Inference Engine (Port 8001))}}\\[0.2cm]
        {\scriptsize สามารถใส่ภาพ \texttt{images/system\_architecture.png} หรือรูปผังระบบที่วาดไว้ลงในจุดนี้}
        \vspace{2.5cm}
    }
}
\caption{แผนภาพสถาปัตยกรรมระบบไมโครเซอร์วิส 3 คอนเทนเนอร์ (Frontend, Backend API Gateway และ AI Inference Engine)}
\label{fig:system_architecture_diagram}
\end{figure}

% ------------------------------------------------------------------------------
% 3.3 การออกแบบผังกรณีการใช้งานและผังกิจกรรมการทำงานของระบบ
% ------------------------------------------------------------------------------
\section{การออกแบบผังกรณีการใช้งานและผังกิจกรรมการทำงานของระบบ}
\label{sec:usecase_and_workflow}

\subsection{ผังกรณีการใช้งานของระบบ (Use Case Diagram)}
การออกแบบผังกรณีการใช้งาน (Use Case Diagram) มีวัตถุประสงค์เพื่อกำหนดขอบเขตหน้าที่และความสามารถของระบบสนับสนุนการวินิจฉัยรอยโรคทางช่องปาก โดยแบ่งกลุ่มผู้มีส่วนร่วมกับระบบ (Actors) ออกเป็น 2 กลุ่มหลัก ได้แก่:
\begin{enumerate}
    \item \textbf{ทันตแพทย์และบุคลากรทางการแพทย์ (Dentist / Medical Personnel):} เป็นผู้ใช้งานหลักของระบบ (Primary Actor) ที่เข้ามาปฏิบัติงานทางคลินิก ได้แก่ การเข้าสู่ระบบ การเลือกสถาปัตยกรรมแบบจำลอง AI ที่ต้องการ การอัปโหลดภาพถ่ายรอยโรค การระบุอาการทางคลินิกและพฤติกรรมเสี่ยงของผู้ป่วย ตลอดจนการตรวจสอบผลการประเมินความเสี่ยงและภาพอธิบายผลทางคลินิก
    \item \textbf{ผู้ดูแลระบบ (System Administrator / Admin):} เป็นกลุ่มผู้ใช้งานที่เตรียมการรองรับสำหรับการขยายขอบเขตระบบในอนาคต (Future Administrative Actor) ทำหน้าที่บริหารจัดการระบบในภาพรวม ได้แก่ การจัดการบัญชีผู้ใช้งานและสิทธิ์การเข้าถึง การจัดการและตรวจสอบสถานะแบบจำลองปัญญาประดิษฐ์ ตลอดจนการตรวจสอบบันทึกประวัติการทำงานของระบบ (System \& Inference Logs)
\end{enumerate}

กรณีการใช้งานหลักของระบบ (Use Cases) จัดหมวดหมู่ตามขอบเขตหน้าที่การทำงานได้ดังนี้:
\begin{enumerate}
    \item \textbf{กลุ่มการยืนยันตัวตน (Authentication):}
    \begin{itemize}
        \item \textbf{UC-AUTH (Login / Authenticate):} การเข้าสู่ระบบเพื่อยืนยันตัวตนสำหรับทันตแพทย์และผู้ดูแลระบบ
    \end{itemize}
    \item \textbf{กลุ่มฟังก์ชันการตรวจคัดกรองทางคลินิกสำหรับทันตแพทย์ (Clinical Screening Functions):}
    \begin{itemize}
        \item \textbf{UC-01 (Select AI Model Architecture):} การเลือกสถาปัตยกรรมแบบจำลองปัญญาประดิษฐ์ที่ต้องการใช้งาน (เช่น แบบจำลองตรวจจับวัตถุ YOLO หรือแบบจำลองหลายโมดาลิตี Multimodal DenseNet-BERT)
        \item \textbf{UC-02 (Upload Lesion Image):} การอัปโหลดภาพถ่ายรอยโรคในช่องปากเข้าสู่ระบบเพื่อเตรียมการประมวลผล
        \item \textbf{UC-03 (Input Patient Symptoms):} การระบุข้อมูลอาการทางคลินิกและประวัติพฤติกรรมเสี่ยงของผู้ป่วย ทั้งการเลือกจากรายการตัวเลือก (Checkboxes) และการกรอกข้อความบรรยายอาการเพิ่มเติม
        \item \textbf{UC-04 (View Clinical Assessment Result):} การเรียกดูผลการประเมินทางคลินิก ภาพผลลัพธ์พร้อมเครื่องมืออธิบายผล (Bounding Box สำหรับ YOLO หรือ Grad-CAM Heatmap สำหรับ Multimodal) ค่าคะแนนความเชื่อมั่น และระดับความเสี่ยงของรอยโรค
    \end{itemize}
    \item \textbf{กลุ่มฟังก์ชันการบริหารจัดการระบบในอนาคตสำหรับผู้ดูแลระบบ (Administrative Functions - Future Extension):}
    \begin{itemize}
        \item \textbf{UC-ADM-01 (Manage User Accounts):} การจัดการข้อมูลบัญชีผู้ใช้งาน กำหนดสิทธิ์ และการอนุมัติการเข้าใช้งานระบบของทันตแพทย์
        \item \textbf{UC-ADM-02 (Manage AI Models):} การจัดการแบบจำลองปัญญาประดิษฐ์ การตรวจสอบสถานะการเชื่อมต่อ และการอัปเดตไฟล์น้ำหนักแบบจำลอง (Model Weights)
        \item \textbf{UC-ADM-03 (View System Logs):} การตรวจสอบบันทึกประวัติการทำงานของระบบ สถิติการเรียกใช้งานการอนุมานผล (Inference Logs) และข้อผิดพลาดที่เกิดขึ้น
    \end{itemize}
\end{enumerate}

ผังกรณีการใช้งานของระบบสนับสนุนการวินิจฉัยรอยโรคทางช่องปาก แสดงความสัมพันธ์ระหว่างตัวแสดง (Actors) และกรณีการใช้งานทั้งหมด ดังแสดงในรูปที่ \ref{fig:usecase_diagram}:

% แผนภาพ Use Case Diagram (เว้นกรอบ Placeholder สำหรับใส่ภาพกราฟิก)
\begin{figure}[htbp]
\centering
\fbox{
    \parbox{0.92\textwidth}{
        \centering
        \vspace{2.5cm}
        \textbf{[พื้นที่สำหรับใส่ภาพ: แผนภาพกรณีการใช้งานของระบบ (Use Case Diagram)]}\\[0.3cm]
        {\footnotesize \textit{(Use Case Diagram: Actor ทันตแพทย์ (Dentist), ผู้ดูแลระบบในอนาคต (Admin) และความสัมพันธ์กับ Use Cases)}}\\[0.2cm]
        {\scriptsize สามารถใส่ภาพ \texttt{images/usecase\_diagram.png} ลงในจุดนี้}
        \vspace{2.5cm}
    }
}
\caption{แผนภาพกรณีการใช้งานของระบบสนับสนุนการวินิจฉัยรอยโรคทางช่องปาก (Use Case Diagram)}
\label{fig:usecase_diagram}
\end{figure}

\subsection{ผังกิจกรรมการทำงานของระบบ (Activity Diagram)}
เพื่ออธิบายลำดับขั้นตอนการประมวลผลคำขอส่งตรวจตั้งแต่จุดเริ่มต้นจนถึงการแสดงผลลัพธ์ ผังกิจกรรมการทำงาน (Activity Diagram) ของระบบแสดงกระบวนการทำงานแบบเป็นขั้นตอน ดังแสดงในรูปที่ \ref{fig:activity_diagram}:

% แผนภาพ Activity Diagram (เว้นกรอบ Placeholder สำหรับใส่ภาพกราฟิก)
\begin{figure}[htbp]
\centering
\fbox{
    \parbox{0.92\textwidth}{
        \centering
        \vspace{2.5cm}
        \textbf{[พื้นที่สำหรับใส่ภาพ: แผนภาพกิจกรรมการทำงานของการประเมินรอยโรค (Activity Diagram)]}\\[0.3cm]
        {\footnotesize \textit{(Activity Diagram: Login $\rightarrow$ Select Model $\rightarrow$ Upload Image $\rightarrow$ Select Symptoms $\rightarrow$ Require NLP? $\rightarrow$ Submit $\rightarrow$ Route to Model $\rightarrow$ View Result)}}\\[0.2cm]
        {\scriptsize สามารถใส่ภาพ \texttt{images/activity\_diagram.png} ลงในจุดนี้}
        \vspace{2.5cm}
    }
}
\caption{แผนภาพกิจกรรมการทำงานของการประเมินรอยโรคและการเลือกแบบจำลอง (Activity Diagram)}
\label{fig:activity_diagram}
\end{figure}

ขั้นตอนการทำงานใน Activity Diagram สามารถอธิบายตามลำดับการไหลของกระบวนการได้ดังนี้:
\begin{enumerate}
    \item \textbf{การเข้าสู่ระบบ (Login):} ผู้ใช้งานทำการเข้าสู่ระบบยืนยันตัวตนเพื่อเข้าถึงหน้าการตรวจคัดกรอง
    \item \textbf{การเลือกแบบจำลอง (Select Model):} ผู้ใช้งานเลือกสถาปัตยกรรมแบบจำลองปัญญาประดิษฐ์ที่ต้องการใช้ในการตรวจรอยโรค
    \item \textbf{การอัปโหลดรูปภาพ (Upload Image):} อัปโหลดภาพถ่ายรอยโรคในช่องปากที่ต้องการประเมินเข้าสู่ระบบ
    \item \textbf{การเลือกอาการทางคลินิก (Select Symptoms Checkbox):} ติ๊กเลือกอาการและพฤติกรรมเสี่ยงจากรายการตัวเลือกที่กำหนด เช่น พฤติกรรมการสูบบุหรี่ ดื่มสุรา เคี้ยวหมาก อาการแสบร้อน หรือลักษณะรอยโรคที่สังเกตได้
    \item \textbf{การตรวจสอบเงื่อนไขความต้องการภาษาธรรมชาติ (Decision Node: Require NLP?):}
    \begin{itemize}
        \item \textbf{กรณีใช่ (YES):} หากแบบจำลองที่เลือกเป็นแบบจำลองในกลุ่มที่ต้องใช้การวิเคราะห์ภาษาธรรมชาติเชิงลึก (เช่น Multimodal WangchanBERTa) ระบบจะแสดงช่องให้ผู้ใช้กรอกข้อความบรรยายอาการและบริบทเพิ่มเติม (\texttt{Fill Symptoms Context}) จากนั้นจึงส่งข้อมูลต่อไปยังขั้นตอนส่งตรวจ
        \item \textbf{กรณีไม่ใช่ (NO):} หากแบบจำลองที่เลือกประมวลผลเฉพาะรูปภาพหรือใช้กฎเกณฑ์จับคู่แบบพื้นฐาน (เช่น YOLO) ผู้ใช้งานสามารถข้ามไปยังขั้นตอนการกดส่งตรวจ (\texttt{Submit}) ได้ทันที
    \end{itemize}
    \item \textbf{การกระจายเส้นทางการประมวลผล (Route based on selection):} Backend API Gateway จะทำการส่งคำขอและข้อมูลไปยัง Inference Pipeline ของแบบจำลองที่ผู้ใช้เลือกไว้ภายใน AI Inference Engine
    \item \textbf{การแสดงผลการประเมิน (View Result):} ระบบรวบรวมผลลัพธ์ คำนวณระดับความเสี่ยง และส่งคืนข้อมูลภาพพร้อมป้ายกำกับเพื่อแสดงผลแก่ผู้ใช้งานบนแดชบอร์ด
\end{enumerate}

% ------------------------------------------------------------------------------
% 3.4 การพัฒนาส่วนติดต่อผู้ใช้งาน
% ------------------------------------------------------------------------------
\section{การพัฒนาส่วนติดต่อผู้ใช้งาน (User Interface Design)}
\label{sec:ui_design}

\subsection{การออกแบบส่วนติดต่อผู้ใช้งานสำหรับการส่งตรวจ (Screening Interface)}
ส่วนติดต่อผู้ใช้ถูกพัฒนาขึ้นโดยเน้นความสะดวกและรวดเร็วในการปฏิบัติงานของทันตแพทย์ ประกอบด้วยฟังก์ชันการทำงานหลัก:
\begin{itemize}
    \item \textbf{ช่องทางเลือกและอัปโหลดภาพ:} รองรับการลากวางไฟล์ภาพ (Drag-and-Drop) พร้อมระบบแสดงตัวอย่างภาพ (Image Preview)
    \item \textbf{เมนูเลือกสถาปัตยกรรมแบบจำลอง (Model Selection Toggle):} ออกแบบเป็นสวิตช์เลือกที่ชัดเจน ให้แพทย์สามารถสลับระหว่างการตรวจจับตำแหน่งด้วย YOLO หรือการวิเคราะห์เชิงลึกด้วย Multimodal DenseNet-BERT ได้อย่างอิสระ
    \item \textbf{ฟอร์มระบุอาการและพฤติกรรมเสี่ยง:} แบ่งหมวดหมู่ชัดเจน ได้แก่ พฤติกรรมสูบบุหรี่ ดื่มสุรา เคี้ยวหมาก อาการแสบร้อนเมื่อกินเผ็ด แผลเจ็บปวด คราบขาวที่เช็ดลอกได้ และตำแหน่งที่พบรอยโรค
\end{itemize}

\subsection{การออกแบบแดชบอร์ดแสดงผลและการเปรียบเทียบผลลัพธ์}
หน้าต่างแสดงผลลัพธ์ทางคลินิก (Clinical Assessment Result View) ได้รับการออกแบบให้แสดงข้อมูลสำคัญครบถ้วนในหน้าจอเดียว:
\begin{enumerate}
    \item \textbf{การแสดงผลภาพเปรียบเทียบ:} แสดงภาพต้นฉบับคู่กับภาพผลลัพธ์ โดยหากเลือกโมเดล YOLO จะแสดงภาพที่มีกรอบ Bounding Box พร้อมค่าความมั่นใจรายกล่อง แต่หากเลือก DenseNet-BERT จะแสดงภาพความร้อน Grad-CAM Heatmap ชี้จุดที่โมเดลใช้ในการตัดสินใจ
    \item \textbf{การแสดงป้ายระดับความเสี่ยง (Risk Level Badge):} แสดงแถบสีเตือนระดับความเสี่ยงตามเกณฑ์ทางคลินิก เช่น สีแดงสำหรับความเสี่ยงสูง (High Risk), สีส้มสำหรับความเสี่ยงปานกลาง (Medium Risk), และสีเขียวสำหรับความเสี่ยงต่ำ (Low Risk)
\end{enumerate}

% ------------------------------------------------------------------------------
% 3.5 การจัดเตรียมสภาพแวดล้อม การคอนเทนเนอร์ไรซ์ และการทดสอบระบบอัตโนมัติ
% ------------------------------------------------------------------------------
\section{การจัดเตรียมสภาพแวดล้อม การคอนเทนเนอร์ไรซ์ และการทดสอบระบบอัตโนมัติ}
\label{sec:devops_and_testing}

\subsection{การจัดสภาพแวดล้อมด้วย Docker และ Docker Compose}
ระบบถูกเขียนไฟล์คอนฟิกูเรชัน \texttt{docker-compose.yml} เพื่อควบคุมการทำงานร่วมกันของทั้ง 3 เซอร์วิสอย่างเป็นเอกภาพ:
\begin{itemize}
    \item \textbf{คอนเทนเนอร์ Frontend (\texttt{oralscan-frontend}):} รันบนสภาพแวดล้อม Node.js ผ่าน Vite Development Server ให้บริการส่วนติดต่อผู้ใช้ที่พอร์ต 5173 โดยเชื่อมต่อไปยัง Backend ผ่านตัวแปรสภาพแวดล้อม \texttt{VITE\_API\_URL}
    \item \textbf{คอนเทนเนอร์ Backend Gateway (\texttt{oralscan-backend}):} ใช้ Base Image \texttt{python:3.11-slim} ติดตั้งแพ็กเกจ FastAPI, Uvicorn, HTTPX และ Pydantic ให้บริการที่พอร์ต 8000 และเชื่อมต่อไปยัง AI Engine ผ่านตัวแปรสภาพแวดล้อม \texttt{AI\_ENGINE\_URL}
    \item \textbf{คอนเทนเนอร์ AI Inference Engine (\texttt{oralscan-ai-engine}):} ใช้ Base Image ที่ติดตั้งไลบรารีสำหรับการคำนวณเชิงลึก ได้แก่ PyTorch, Torchvision, Transformers, OpenCV-headless, และ Ultralytics ให้บริการที่พอร์ต 8001
    \item \textbf{การเชื่อมต่อเครือข่ายภายใน (Bridge Network - \texttt{oralscan-network}):} ทั้งสามคอนเทนเนอร์เชื่อมต่อสื่อสารกันผ่านวงเครือข่ายเสมือนภายใน ช่วยเพิ่มความปลอดภัยเนื่องจาก AI Engine ไม่จำเป็นต้องเปิดพอร์ตออกสู่เครือข่ายอินเทอร์เน็ตภายนอกโดยตรง
\end{itemize}

\subsection{การทดสอบระบบอัตโนมัติ (Automated API Testing)}
คณะผู้จัดทำได้เขียนสคริปต์ทดสอบระบบอัตโนมัติ (\texttt{test\_api.py}) เพื่อตรวจสอบความถูกต้องของระบบก่อนการส่งมอบ ครอบคลุมกรณีทดสอบหลักดังแสดงในตารางที่ \ref{tab:test_cases}:

\begin{table}[htbp]
\centering
\small
\caption{ชุดกรณีทดสอบอัตโนมัติ (Automated Test Cases) สำหรับตรวจสอบการทำงานของระบบ}
\label{tab:test_cases}
\begin{tabular}{p{2.2cm}p{4.2cm}p{3.5cm}p{3.5cm}}
\toprule
\textbf{รหัสทดสอบ} & \textbf{เงื่อนไขและข้อมูลนำเข้า} & \textbf{ผลลัพธ์ที่คาดหวัง} & \textbf{เกณฑ์การผ่าน} \\
\midrule
TC-SYS-01 & ส่งคำขอ \texttt{GET /api/health} & สถานะ HTTP 200 ระบบพร้อมทำงาน & ได้รับ JSON \texttt{status: healthy} \\
TC-SYS-02 & ส่งภาพและอาการด้วยโมเดล \texttt{YOLO} & ได้รับผลการทำนายและ Bounding Box & ได้รับ Bounding Box ถูกต้องตาม Schema \\
TC-SYS-03 & ส่งภาพและอาการด้วยโมเดล \texttt{DenseNet-BERT} & ได้รับผลความเสี่ยงและภาพ Grad-CAM & ได้รับผลการวิเคราะห์และ Heatmap URL \\
TC-SYS-04 & ส่งไฟล์ที่ไม่ใช่ภาพถ่าย (เช่น .pdf, .txt) & สถานะ HTTP 400 Bad Request & ปฏิเสธคำขอและแจ้งเตือนรูปแบบไฟล์ \\
TC-SYS-05 & ส่งคำขอโดยระบุ \texttt{modelType} ผิดพลาด & สถานะ HTTP 400 Bad Request & ปฏิเสธคำขอและแนะนำโมเดลที่ถูกต้อง \\
TC-SYS-06 & ทดสอบการส่งคำขอพร้อมกัน 20 คำขอ & ทุกคำขอประมวลผลสำเร็จโดยไม่เกิด Timeout & อัตราความสำเร็จร้อยละ 100 \\
\bottomrule
\end{tabular}
\end{table}

% ------------------------------------------------------------------------------
% 3.6 สรุปวิธีการดำเนินงาน
% ------------------------------------------------------------------------------
\section{สรุปวิธีการดำเนินงาน}
\label{sec:methodology_summary}

การดำเนินงานในบทที่ 3 ได้มุ่งเน้นการแก้ปัญหาข้อจำกัดของงานวิจัยต้นแบบในอดีต โดยนำแบบจำลองตรวจจับวัตถุ YOLO จากรุ่นพี่กลุ่มที่ 1 และแบบจำลองการเรียนรู้เชิงลึกแบบหลายโมดาลิตีจากรุ่นพี่กลุ่มที่ 2 มารวบรวมและบูรณาการเข้าด้วยกันภายใต้สถาปัตยกรรมไมโครเซอร์วิส 3 คอนเทนเนอร์ (Frontend, API Gateway และ AI Engine) ที่ทันสมัยและประหยัดทรัพยากร พร้อมทั้งสร้างทางเลือกให้ทันตแพทย์สามารถเลือกใช้แบบจำลองได้ตามความต้องการของเคสผู้ป่วย (Alternative Model Selection) ผ่านการวางผังกรณีการใช้งาน (Use Case Diagram) และผังกิจกรรมการประมวลผลคำขอ (Activity Diagram) 

ผลการทดสอบการทำงานของระบบ ประสิทธิภาพความเร็วในการประมวลผล และผลการเปรียบเทียบความแม่นยำของแบบจำลองแต่ละแบบบนแพลตฟอร์มนี้ จะถูกนำเสนออย่างละเอียดในบทที่ 4 ต่อไป

    \chapter{ผลการทดลองเบื้องต้นหรือระบบต้นแบบ}
\label{chapter:result}

บทที่สี่

    \chapter{บทสรุป}
\label{chapter:conclusion}

บทที่ห้า


    \clearpage

    \def\inbib{1}
    \addcontentsline{toc}{chapter}{\bibname}
    \bibliographystyle{IEEEtran-kmitl}
    \makeatletter
    \EveryShipout{
        \ifdim\pagetotal>\pagegoal
            \if\inbib1
                \aftergroup\@cont@heading
            \else
            \fi
        \fi
    }
    \makeatother

    \bibliography{reference}
    \def\inbib{0}

    \startappendix
    % ==============================================================================
% ภาคผนวก (Appendix)
% โครงการ: การพัฒนาแพลตฟอร์มสำหรับการสนับสนุนการวินิจฉัยรอยโรคทางช่องปาก
%          และผสมผสานการใช้งานระบบผู้เชี่ยวชาญสำหรับทันตแพทย์
% ==============================================================================

% ------------------------------------------------------------------------------
% ภาคผนวก ก: คู่มือการติดตั้งและกำหนดค่าสภาพแวดล้อมระบบ
% ------------------------------------------------------------------------------
\chapter{คู่มือการติดตั้งและการกำหนดค่าสภาพแวดล้อมระบบ}
\label{appendix:installation}

ภาคผนวกนี้อธิบายข้อกำหนดของระบบ โครงสร้างไฟล์ และขั้นตอนการติดตั้งสภาพแวดล้อมการทำงานของระบบสนับสนุนการวินิจฉัยรอยโรคทางช่องปาก โดยใช้เทคโนโลยีคอนเทนเนอร์ (Docker และ Docker Compose) เพื่ออำนวยความสะดวกในการนำระบบไปติดตั้งและใช้งานจริง

\section{ข้อกำหนดของระบบ (System Requirements)}
เพื่อให้ระบบสามารถประมวลผลแบบจำลองปัญญาประดิษฐ์และให้บริการเว็บแอปพลิเคชันได้อย่างมีประสิทธิภาพ อุปกรณ์เซิร์ฟเวอร์ควรมีคุณสมบัติขั้นต่ำดังนี้:

\begin{enumerate}
    \item \textbf{ข้อกำหนดด้านฮาร์ดแวร์ (Hardware Requirements):}
    \begin{itemize}
        \item หน่วยประมวลผลกลาง (CPU): สถาปัตยกรรม x86\_64 ความเร็วไม่น้อยกว่า 2.0 GHz จำนวน 4 แกนประมวลผล (Cores) ขึ้นไป
        \item หน่วยความจำหลัก (RAM): ไม่น้อยกว่า 8 GB (แนะนำ 16 GB ขึ้นไป สำหรับการรันแบบจำลอง Multimodal คู่ขนาน)
        \item พื้นที่จัดเก็บข้อมูล (Storage): ไม่น้อยกว่า 15 GB (สำหรับจัดเก็บ Docker Images, Weights ของโมเดล และไฟล์ภาพประวัติ)
        \item หน่วยประมวลผลกราฟิก (GPU - ทางเลือก): NVIDIA GPU รองรับ CUDA 11.8+ พร้อม VRAM ไม่น้อยกว่า 6 GB (หากไม่มี ระบบจะสลับไปประมวลผลบน CPU อัตโนมัติ)
    \end{itemize}
    \item \textbf{ข้อกำหนดด้านซอฟต์แวร์ (Software Requirements):}
    \begin{itemize}
        \item ระบบปฏิบัติการ: Ubuntu Linux 22.04 LTS, Debian 12 หรือ Windows 11 ที่เปิดใช้งาน WSL2 (Windows Subsystem for Linux)
        \item ซอฟต์แวร์คอนเทนเนอร์: Docker Engine เวอร์ชัน 24.0.0 ขึ้นไป และ Docker Compose Plugin เวอร์ชัน 2.20.0 ขึ้นไป
        \item เว็บบราวเซอร์สำหรับผู้ใช้งาน: Google Chrome, Mozilla Firefox หรือ Microsoft Edge เวอร์ชันล่าสุด
    \end{itemize}
\end{enumerate}

\section{โครงสร้างไดเรกทอรีของโครงงาน (Project Directory Structure)}
โครงสร้างแฟ้มข้อมูลหลักของระบบสนับสนุนการวินิจฉัยรอยโรคทางช่องปาก ถูกจัดระเบียบตามสถาปัตยกรรมไมโครเซอร์วิส ดังนี้:

\begin{verbatim}
Oral_Lesion_2026/
|-- docker-compose.yml              # ไฟล์คอนฟิกูเรชันสำหรับควบคุมทุกคอนเทนเนอร์
|-- test_api.py                     # สคริปต์ทดสอบการทำงานของ API อัตโนมัติ
|-- frontend/                       # ไมโครเซอร์วิสส่วนติดต่อผู้ใช้งาน (React + Vite)
|   |-- Dockerfile
|   |-- package.json
|   |-- src/
|   `-- tailwind.config.js
|-- backend/                        # ไมโครเซอร์วิส Backend API Gateway (FastAPI)
|   |-- Dockerfile
|   |-- requirements.txt
|   |-- main.py
|   `-- static/
|       |-- uploads/                # ไดเรกทอรีเก็บภาพที่อัปโหลดเข้าสู่ระบบ
|       `-- outputs/                # ไดเรกทอรีเก็บภาพผลลัพธ์ (BBox / Grad-CAM)
`-- ai-engine/                      # ไมโครเซอร์วิสประมวลผลปัญญาประดิษฐ์
    |-- Dockerfile
    |-- requirements.txt
    |-- main.py
    |-- yolo-service/               # โมดูลตรวจจับวัตถุ YOLOv7 และ YOLOv8
    |   |-- yolov7.pt
    |   `-- yolov8.pt
    `-- densenet-bert-service/      # โมดูลการเรียนรู้เชิงลึกแบบหลายโมดาลิตี
        |-- best_densenet121_img.pth
        `-- best_text.pth
\end{verbatim}

\section{การจัดเตรียมไฟล์น้ำหนักแบบจำลอง (Model Weights Setup)}
ก่อนเริ่มต้นการทำงานของระบบ ต้องตรวจสอบให้แน่ใจว่าได้วางไฟล์น้ำหนักแบบจำลอง (Pre-trained Weights) ไว้ในตำแหน่งที่ถูกต้อง ดังนี้:
\begin{enumerate}
    \item ไฟล์แบบจำลอง YOLOv7: จัดเก็บที่ \texttt{ai-engine/yolo-service/yolov7.pt}
    \item ไฟล์แบบจำลอง YOLOv8: จัดเก็บที่ \texttt{ai-engine/yolo-service/yolov8.pt}
    \item ไฟล์แบบจำลอง DenseNet-121 สำหรับประมวลผลภาพ: จัดเก็บที่ \texttt{ai-engine/densenet-bert-service/best\_densenet121\_img.pth}
    \item ไฟล์แบบจำลอง WangchanBERTa สำหรับประมวลผลข้อความ: จัดเก็บที่ \texttt{ai-engine/densenet-bert-service/best\_text.pth}
\end{enumerate}

\section{การกำหนดค่าตัวแปรสภาพแวดล้อม (Environment Variables)}
ระบบควบคุมการทำงานของแต่ละไมโครเซอร์วิสผ่านตัวแปรสภาพแวดล้อมในไฟล์ \texttt{docker-compose.yml} ดังแสดงในตารางที่ \ref{tab:env_vars}:

\begin{table}[htbp]
\centering
\small
\caption{รายละเอียดตัวแปรสภาพแวดล้อมของระบบ}
\label{tab:env_vars}
\begin{tabular}{p{2.5cm}p{3.5cm}p{3.5cm}p{4cm}}
\toprule
\textbf{เซอร์วิส} & \textbf{ชื่อตัวแปร} & \textbf{ค่าเริ่มต้น} & \textbf{คำอธิบายหน้าที่} \\
\midrule
Frontend & \texttt{VITE\_API\_URL} & \texttt{http://localhost:8000} & ที่อยู่ URL สำหรับเชื่อมต่อ Backend Gateway \\
\midrule
Backend & \texttt{AI\_ENGINE\_URL} & \texttt{http://ai-engine:8001} & ที่อยู่ภายในวงเครือข่ายสำหรับเรียก AI Engine \\
\midrule
\multirow{2}{*}{AI Engine} & \texttt{PYTHONUNBUFFERED} & \texttt{1} & บังคับให้ Python แสดงผล log ออกทันที \\
 & \texttt{FUSION\_ALPHA} & \texttt{0.6} & ค่าน้ำหนักถ่วงภาพในการคำนวณ Late Fusion \\
\bottomrule
\end{tabular}
\end{table}

\section{ขั้นตอนการสร้างและเริ่มต้นระบบ (Deployment Steps)}
ผู้ดูแลระบบสามารถสั่งสร้างและเปิดให้บริการระบบทั้งหมดผ่านคำสั่ง Docker Compose ตามขั้นตอนดังนี้:

\begin{enumerate}
    \item เปิดเทอร์มินัลและเข้าไปยังไดเรกทอรีรากของโครงงาน:
\begin{verbatim}
cd Oral_Lesion_2026
\end{verbatim}

    \item สั่งสร้างคอนเทนเนอร์อิมเมจและเริ่มต้นระบบในโหมดเบื้องหลัง (Detached Mode):
\begin{verbatim}
docker compose up -d --build
\end{verbatim}

    \item ตรวจสอบสถานะความพร้อมของคอนเทนเนอร์ทั้ง 3 เซอร์วิส:
\begin{verbatim}
docker compose ps
\end{verbatim}
    เมื่อระบบพร้อมทำงาน คอนเทนเนอร์ \texttt{oralscan-frontend}, \texttt{oralscan-backend}, และ \texttt{oralscan-ai-engine} จะต้องแสดงสถานะเป็น \texttt{Up (healthy)} หรือ \texttt{Up}

    \item การเข้าใช้งานระบบผ่านเว็บบราวเซอร์:
    \begin{itemize}
        \item ส่วนติดต่อผู้ใช้งานสำหรับทันตแพทย์: \url{http://localhost:5173}
        \item เอกสาร API ทางการ (Swagger UI): \url{http://localhost:8000/docs}
    \end{itemize}

    \item คำสั่งสำหรับหยุดและปิดการทำงานของระบบ:
\begin{verbatim}
docker compose down
\end{verbatim}
\end{enumerate}

% ------------------------------------------------------------------------------
% ภาคผนวก ข: ข้อกำหนดส่วนต่อประสานโปรแกรมประยุกต์
% ------------------------------------------------------------------------------
\chapter{ข้อกำหนดส่วนต่อประสานโปรแกรมประยุกต์ (API Specification)}
\label{appendix:api_spec}

ภาคผนวกนี้อธิบายข้อกำหนดทางเทคนิคของ RESTful API ที่ให้บริการโดย Backend API Gateway สำหรับการเชื่อมต่อระหว่างระบบส่วนหน้า (Frontend) หรือระบบภายนอกกับระบบประเมินรอยโรคช่องปาก

\section{ภาพรวมของส่วนต่อประสาน (API Overview)}
ระบบให้บริการ API ภายใต้โปรโตคอล HTTP/1.1 และส่งคืนข้อมูลในรูปแบบ JSON (JavaScript Object Notation) โดยมี Endpoint หลักแสดงในตารางที่ \ref{tab:api_endpoints}:

\begin{table}[htbp]
\centering
\small
\caption{รายการ API Endpoints ของระบบสนับสนุนการวินิจฉัยรอยโรคทางช่องปาก}
\label{tab:api_endpoints}
\begin{tabular}{p{1.8cm}p{4.2cm}p{3.5cm}p{4.2cm}}
\toprule
\textbf{Method} & \textbf{Endpoint Path} & \textbf{Content-Type} & \textbf{คำอธิบายหน้าที่} \\
\midrule
\texttt{GET} & \texttt{/api/health} & \texttt{application/json} & ตรวจสอบสถานะความพร้อมของเซิร์ฟเวอร์ \\
\texttt{POST} & \texttt{/api/assess} & \texttt{multipart/form-data} & ส่งภาพถ่ายและข้อมูลอาการเพื่อประเมินรอยโรค \\
\texttt{GET} & \texttt{/static/outputs/\{file\}} & \texttt{image/jpeg} หรือ \texttt{image/png} & ให้บริการเรียกดูไฟล์ภาพผลลัพธ์ \\
\bottomrule
\end{tabular}
\end{table}

\section{ข้อกำหนดของ Endpoint ประเมินรอยโรค (/api/assess)}

\subsection{ข้อมูลนำเข้า (Request Parameters)}
คำขอส่งตรวจประเมินรอยโรคต้องส่งด้วยวิธี \texttt{POST} ในรูปแบบ \texttt{multipart/form-data} ประกอบด้วยฟิลด์ข้อมูลดังนี้:

\begin{enumerate}
    \item \textbf{\texttt{image} (จำเป็น):} ไฟล์ภาพถ่ายรอยโรคในช่องปาก รองรับนามสกุลไฟล์ \texttt{.jpg}, \texttt{.jpeg}, และ \texttt{.png} ขนาดไฟล์ไม่เกิน 10 MB
    \item \textbf{\texttt{modelType} (จำเป็น):} ข้อความระบุชื่อสถาปัตยกรรมแบบจำลองที่ต้องการใช้ประมวลผล รองรับ 3 ค่า ได้แก่:
    \begin{itemize}
        \item \texttt{"YOLOv7"} สำหรับตรวจจับตำแหน่งด้วยแบบจำลอง YOLOv7
        \item \texttt{"YOLOv8"} สำหรับตรวจจับตำแหน่งด้วยแบบจำลอง YOLOv8
        \item \texttt{"DenseNet-BERT"} สำหรับการจำแนกหลายโมดาลิตีและสร้างแผนที่ Grad-CAM Heatmap
    \end{itemize}
    \item \textbf{\texttt{text} (ทางเลือก):} ข้อความบรรยายอาการ พฤติกรรมเสี่ยง หรือคำสำคัญทางคลินิก (เช่น \texttt{"smoking, ulcer, ดื่มเหล้า"}) ซึ่งจำเป็นต้องระบุหากเลือกใช้งานโมเดล \texttt{"DenseNet-BERT"} เพื่อให้โมดูล WangchanBERTa นำไปประมวลผล
\end{enumerate}

\subsection{โครงสร้างข้อมูลตอบกลับสำเร็จ (Success Response: HTTP 200 OK)}
เมื่อระบบประมวลผลคำขอสำเร็จ จะส่งคืนข้อมูล JSON ตามโครงสร้างตัวอย่างดังนี้:

\begin{verbatim}
{
  "success": true,
  "model_used": "YOLOv7",
  "result": {
    "disease_name": "Leukoplakia",
    "risk_level": "High",
    "confidence_score": 0.884,
    "bounding_boxes": [
      {
        "xmin": 142.5,
        "ymin": 210.0,
        "xmax": 358.2,
        "ymax": 420.7,
        "label": "Leukoplakia",
        "conf": 0.884
      }
    ],
    "processed_image_url": "/static/outputs/annotated_uuid123.jpg"
  }
}
\end{verbatim}

\subsection{รหัสสถานะและการจัดการข้อผิดพลาด (Error Handling)}
ระบบมีการตรวจสอบความถูกต้องของข้อมูลนำเข้าและส่งคืนรหัสสถานะ HTTP ตามมาตรฐานสากล:
\begin{itemize}
    \item \textbf{\texttt{400 Bad Request}:} เกิดขึ้นเมื่อข้อมูลนำเข้าไม่ถูกต้อง เช่น ไม่อัปโหลดไฟล์ภาพ, นามสกุลไฟล์ไม่ใช่ภาพ, หรือระบุชื่อ \texttt{modelType} ไม่ตรงกับที่ระบบรองรับ
    \item \textbf{\texttt{500 Internal Server Error}:} เกิดขึ้นเมื่อโมดูล AI Engine ขัดข้อง หรือหน่วยความจำไม่เพียงพอขณะทำการอนุมานผล
\end{itemize}

% ------------------------------------------------------------------------------
% ภาคผนวก ค: สคริปต์การทดสอบและผลการทดสอบระบบอัตโนมัติ
% ------------------------------------------------------------------------------
\chapter{สคริปต์การทดสอบและผลการทดสอบระบบอัตโนมัติ}
\label{appendix:test_scripts}

ภาคผนวกนี้แสดงซอร์สโค้ดของสคริปต์การทดสอบระบบอัตโนมัติ (\texttt{test\_api.py}) ซึ่งใช้สำหรับตรวจสอบการทำงานของ Backend API และตัวอย่างผลลัพธ์การตอบกลับที่ได้รับจากการทดสอบจริง

\section{ซอร์สโค้ดสคริปต์การทดสอบ (test\_api.py)}
สคริปต์ทดสอบถูกพัฒนาขึ้นด้วยภาษา Python โดยใช้ไลบรารีมาตรฐาน \texttt{urllib.request} เพื่อจำลองการส่งคำขอแบบ \texttt{multipart/form-data} ไปยัง Endpoint \texttt{/api/assess}:

\begin{verbatim}
import urllib.request
import json
from pathlib import Path

def test_endpoint(model_type, text=""):
    url = "http://localhost:8000/api/assess"
    boundary = "----WebKitFormBoundaryTest123"
    
    parts = []
    # 1. ฟิลด์ระบุประเภทโมเดล (modelType)
    parts.append(
        f"--{boundary}\r\n"
        f"Content-Disposition: form-data; name=\"modelType\"\r\n\r\n"
        f"{model_type}\r\n".encode("utf-8")
    )
    
    # 2. ฟิลด์ระบุข้อความอาการ (text)
    if text:
        parts.append(
            f"--{boundary}\r\n"
            f"Content-Disposition: form-data; name=\"text\"\r\n\r\n"
            f"{text}\r\n".encode("utf-8")
        )
        
    # 3. แนบไฟล์ภาพตัวอย่าง (image)
    parts.append(
        f"--{boundary}\r\n"
        f"Content-Disposition: form-data; name=\"image\"; filename=\"test.jpg\"\r\n"
        f"Content-Type: image/jpeg\r\n\r\n".encode("utf-8")
    )
    with open("static/uploads/test.jpg", "rb") as f:
        parts.append(f.read())
    parts.append(f"\r\n--{boundary}--\r\n".encode("utf-8"))
    
    body = b"".join(parts)
    req = urllib.request.Request(
        url,
        data=body,
        headers={"Content-Type": f"multipart/form-data; boundary={boundary}"}
    )
    
    with urllib.request.urlopen(req) as resp:
        res = json.loads(resp.read().decode("utf-8"))
        print(f"\n[+] Test {model_type} -> HTTP {resp.status}")
        print(json.dumps(res, indent=2, ensure_ascii=False))

if __name__ == "__main__":
    print("Testing YOLOv7 Model...")
    test_endpoint("YOLOv7", "smoking")
    
    print("\nTesting YOLOv8 Model...")
    test_endpoint("YOLOv8", "alwaysHurts")
    
    print("\nTesting DenseNet-BERT Model...")
    test_endpoint("DenseNet-BERT", "ดื่มเครื่องดื่มแอลกอฮอล์ สูบบุหรี่ ปื้นขาว ลิ้น")
\end{verbatim}

\section{ตัวอย่างผลลัพธ์จากการทดสอบแบบจำลองแต่ละประเภท}

\subsection{ผลการทดสอบแบบจำลอง YOLOv7}
ผลลัพธ์การตอบกลับเมื่อส่งทดสอบภาพถ่ายรอยโรคด้วยโมเดล YOLOv7 พร้อมคำสำคัญ \texttt{"smoking"}:

\begin{verbatim}
[+] Test YOLOv7 -> HTTP 200
{
  "success": true,
  "model_used": "YOLOv7",
  "result": {
    "disease_name": "Leukoplakia",
    "risk_level": "High",
    "confidence_score": 0.884,
    "bounding_boxes": [
      {
        "xmin": 154.2,
        "ymin": 204.6,
        "xmax": 372.8,
        "ymax": 418.1,
        "label": "Leukoplakia",
        "conf": 0.884
      }
    ],
    "processed_image_url": "/static/outputs/annotated_yolo7_f8a2c1.jpg"
  }
}
\end{verbatim}

\subsection{ผลการทดสอบแบบจำลอง Multimodal DenseNet-BERT}
ผลลัพธ์การตอบกลับเมื่อส่งทดสอบภาพถ่ายรอยโรคด้วยโมเดล Multimodal DenseNet-BERT ร่วมกับข้อความอาการภาษาไทย \texttt{"ดื่มเครื่องดื่มแอลกอฮอล์ สูบบุหรี่ ปื้นขาว ลิ้น"}:

\begin{verbatim}
[+] Test DenseNet-BERT -> HTTP 200
{
  "success": true,
  "model_used": "DenseNet-BERT",
  "result": {
    "disease_name": "Leukoplakia",
    "risk_level": "High",
    "confidence_score": 0.912,
    "bounding_boxes": [],
    "processed_image_url": "/static/outputs/gradcam_densenet_3b91ef.jpg"
  }
}
\end{verbatim}

% ------------------------------------------------------------------------------
% ภาคผนวก ง: พจนานุกรมอาการทางคลินิกและเกณฑ์การประเมินความเสี่ยง
% ------------------------------------------------------------------------------
\chapter{พจนานุกรมอาการทางคลินิกและเกณฑ์การประเมินความเสี่ยง}
\label{appendix:clinical_dictionary}

ภาคผนวกนี้แสดงรายละเอียดชุดคำสำคัญทางคลินิกสำหรับใช้ในการจับคู่กฎเกณฑ์ (Rule Matching) และเกณฑ์การจำแนกระดับความเสี่ยงของรอยโรคที่นำมาใช้ในระบบสนับสนุนการตัดสินใจทางคลินิก

\section{พจนานุกรมคำสำคัญทางคลินิกจำเพาะตามรอยโรค}
ตารางที่ \ref{tab:clinical_symptoms} แสดงรายการคำสำคัญและอาการที่ใช้ตรวจสอบความสอดคล้องระหว่างข้อมูลซักประวัติของผู้ป่วยกับชนิดรอยโรคที่แบบจำลองตรวจพบ:

\begin{table}[htbp]
\centering
\small
\caption{พจนานุกรมคำสำคัญและอาการทางคลินิกจำเพาะตามรอยโรคช่องปาก}
\label{tab:clinical_symptoms}
\begin{tabular}{p{3.2cm}p{5.5cm}p{4.5cm}}
\toprule
\textbf{ชื่อรอยโรค (Pathology)} & \textbf{คำสำคัญภาษาไทยที่ระบบรองรับ} & \textbf{รหัสฟิลด์ข้อมูลในระบบ} \\
\midrule
Leukoplakia & ดื่มเหล้า, สูบบุหรี่, เคี้ยวหมาก, แอลกอฮอล์ & \texttt{drinkalcohol}, \texttt{smoking}, \texttt{chewbetelnut} \\
\midrule
Lichen planus & กินเผ็ด, กินเผ็ดแสบ, รับประทานอาหารเผ็ด & \texttt{eatspicyfood} \\
\midrule
Candidiasis & เช็ดออก, เช็ดออกได้, คราบขาวที่สามารถลอกออก & \texttt{wipeoff} \\
\midrule
Linea alba buccalis & ไม่มีอาการ, ไม่มีอาการผิดปกติ & \texttt{nosymptoms} \\
\midrule
Ulcer & เจ็บเมื่อโดนแผล, มีอาการเจ็บ, เจ็บตลอดเวลา, ระคายเคืองตลอดเวลา & \texttt{alwayshurts} \\
\bottomrule
\end{tabular}
\end{table}

\section{เกณฑ์การจัดระดับความเสี่ยงของรอยโรคทางคลินิก (Risk Stratification)}
ระบบจัดระดับความเสี่ยง (Risk Level) เพื่อช่วยให้ทันตแพทย์สามารถจัดลำดับความเร่งด่วนในการส่งต่อหรือตรวจชิ้นเนื้อเพิ่มเติม (Biopsy) โดยแบ่งออกเป็น 3 ระดับตามตารางที่ \ref{tab:risk_levels}:

\begin{table}[htbp]
\centering
\small
\caption{เกณฑ์การจำแนกระดับความเสี่ยงของรอยโรคทางช่องปาก}
\label{tab:risk_levels}
\begin{tabular}{p{2.5cm}p{4cm}p{3cm}p{4cm}}
\toprule
\textbf{ระดับความเสี่ยง} & \textbf{ชนิดรอยโรคที่เข้าเกณฑ์} & \textbf{เงื่อนไขความมั่นใจ} & \textbf{คำแนะนำในการปฏิบัติทางคลินิก} \\
\midrule
\textbf{ความเสี่ยงสูง\newline(High Risk)} & Leukoplakia, Ulcer, มะเร็งช่องปาก (Oral Cancer) & ค่าคะแนนความเชื่อมั่น $\ge 0.50$ & มีโอกาสพัฒนาเป็นมะเร็งช่องปาก แนะนำให้ตรวจชิ้นเนื้อหรือส่งต่อแพทย์ผู้เชี่ยวชาญ \\
\midrule
\textbf{ความเสี่ยงปานกลาง\newline(Medium Risk)} & Lichen planus, หรือกลุ่มความเสี่ยงสูงที่มีความเชื่อมั่นต่ำกว่าเกณฑ์ & ค่าคะแนนความเชื่อมั่น $\ge 0.60$ & แนะนำให้นัดหมายติดตามอาการอย่างใกล้ชิดภายใน 2-4 สัปดาห์ \\
\midrule
\textbf{ความเสี่ยงต่ำ\newline(Low Risk)} & Candidiasis, Linea alba buccalis, หรือรอยโรคทั่วไป & ทุกระดับความเชื่อมั่น & รอยโรคไม่ใช่มะเร็ง รักษาตามอาการหรือให้คำแนะนำการรักษาสุขอนามัยช่องปาก \\
\bottomrule
\end{tabular}
\end{table}

\section{การกำหนดสีสัญลักษณ์และรูปแบบการแสดงผลทางคลินิก}
เพื่อช่วยให้ทันตแพทย์สามารถสังเกตความผิดปกติได้อย่างรวดเร็ว ระบบได้กำหนดสีของเส้นกรอบ Bounding Box จำเพาะตามประเภทของรอยโรค ดังนี้:
\begin{itemize}
    \item \textbf{Leukoplakia:} เส้นกรอบสีเขียว (Green - \texttt{BGR: 0, 255, 0})
    \item \textbf{Lichen planus:} เส้นกรอบสีแดง (Red - \texttt{BGR: 0, 0, 255})
    \item \textbf{Candidiasis:} เส้นกรอบสีขาว (White - \texttt{BGR: 255, 255, 255})
    \item \textbf{Linea alba buccalis:} เส้นกรอบสีเทา (Gray - \texttt{BGR: 128, 128, 128})
    \item \textbf{Ulcer:} เส้นกรอบสีส้ม (Orange - \texttt{BGR: 255, 140, 0})
\end{itemize}


    \makeauthorbio

\end{document}